
\Data\ID{1003099803}
\Updated{2022-04-23 05:04:57}
\Level{A}
\SubArea{Logarithmic Functions}
\LANGen{
\Question{
Find the vertical asymptote of \( f(x)=\log_{\frac13}(x-2) \).}
{\( x =2 \)}{\( x =-2 \)}{\( y =2 \)}{\( y =-2 \)}{}{}}
\LANGpl{
\Question{
Znajdź pionową asymptotę funkcji \( f(x)=\log_{\frac13}(x-2) \).}
{\( x =2 \)}{\( x =-2 \)}{\( y =2 \)}{\( y =-2 \)}{}{}}
\LANGcs{
\Question{
Která rovnice je rovnicí svislé asymptoty funkce \( f(x)=\log_{\frac13}(x-2) \)?}
{\( x =2 \)}{\( x =-2 \)}{\( y =2 \)}{\( y =-2 \)}{}{}}
\LANGsk{
\Question{
Ktorá rovnica je rovnicou zvislej asymptoty funkcie \( f(x)=\log_{\frac13}(x-2) \)?}
{\( x =2 \)}{\( x =-2 \)}{\( y =2 \)}{\( y =-2 \)}{}{}}
\LANGes{
\Question{
Halla la asíntota vertical de la función \( f(x)=\log_{\frac13}(x-2) \).}
{\( x =2 \)}{\( x =-2 \)}{\( y =2 \)}{\( y =-2 \)}{}{}}
\End

\Data\ID{1003099804}
\Updated{2021-12-16 11:12:31}
\Level{A}
\SubArea{Logarithmic Functions}
\LANGen{
\Question{
Which of the following equations is the equation of the vertical asymptote of  \( f(x)=\log_3(x+5) \)?}
{\( x=-5 \)}{\( x=5 \)}{\( y=5 \)}{\( y=-5 \)}{}{}}
\LANGpl{
\Question{
Które z podanych równań jest równaniem pionowej asymptoty  \( f(x)=\log_3(x+5) \)?}
{\( x=-5 \)}{\( x=5 \)}{\( y=5 \)}{\( y=-5 \)}{}{}}
\LANGcs{
\Question{
Která rovnice je rovnicí svislé asymptoty funkce \( f(x)=\log_3(x+5) \)?}
{\( x=-5 \)}{\( x=5 \)}{\( y=5 \)}{\( y=-5 \)}{}{}}
\LANGsk{
\Question{
Ktorá rovnica je rovnicou zvislej asymptoty funkcie \( f(x)=\log_3(x+5) \)?}
{\( x=-5 \)}{\( x=5 \)}{\( y=5 \)}{\( y=-5 \)}{}{}}
\LANGes{
\Question{
¿Cuál de las siguientes ecuaciones es la ecuación de la asíntota vertical de la función  \( f(x)=\log_3(x+5) \)?}
{\( x=-5 \)}{\( x=5 \)}{\( y=5 \)}{\( y=-5 \)}{}{}}
\End

\Data\ID{1003100001}
\Updated{2021-04-12 08:04:03}
\Level{A}
\SubArea{Logarithmic Functions}
\LANGen{
\Question{
In which of the following options is not given a logarithmic function?}
{\( f(x) = \log_{-2}x\)}{\( g(x) = \log_{2}x\)}{\( h(x) = \log_{\frac12}x\)}{\( m(x) = \log_{0.2}x\)}{}{}}
\LANGpl{
\Question{
W której z poniższych opcji nie jest dana funkcja logarytmiczna?}
{\( f(x) = \log_{-2}x\)}{\( g(x) = \log_{2}x\)}{\( h(x) = \log_{\frac12}x\)}{\( m(x) = \log_{0{,}2}x\)}{}{}}
\LANGcs{
\Question{
V které z následujících možností není uvedená logaritmická funkce?}
{\( f(x) = \log_{-2}x\)}{\( g(x) = \log_{2}x\)}{\( h(x) = \log_{\frac12}x\)}{\( m(x) = \log_{0{,}2}x\)}{}{}}
\LANGsk{
\Question{
V ktorej z nasledujúcich možností nie je uvedená logaritmická funkcia?}
{\( f(x) = \log_{-2}x\)}{\( g(x) = \log_{2}x\)}{\( h(x) = \log_{\frac12}x\)}{\( m(x) = \log_{0{,}2}x\)}{}{}}
\LANGes{
\Question{
¿Qué opción no puede ser una función logarítmica?}
{\( f(x) = \log_{-2}x\)}{\( g(x) = \log_{2}x\)}{\( h(x) = \log_{\frac12}x\)}{\( m(x) = \log_{0.2}x\)}{}{}}
\End

\Data\ID{1103099701}
\Updated{2021-04-12 08:04:53}
\Level{A}
\SubArea{Logarithmic Functions}
\LANGen{
\Question{
One of the given graphs is the graph of a logarithmic function. Choose this graph.}
{}{}{}{}{}{}}
\LANGpl{
\Question{
Który z  poniższych wykresów przedstawia funkcję logarytmiczną?}
{}{}{}{}{}{}}
\LANGcs{
\Question{
Který z následujících grafů je grafem logaritmické funkce?}
{}{}{}{}{}{}}
\LANGsk{
\Question{
Ktorý z nasledujúcich grafov je graf logaritmickej funkcie?}
{}{}{}{}{}{}}
\LANGes{
\Question{
De las gráficas dadas, elige la de una función logarítmica.}
{}{}{}{}{}{}}

\IMGanswer{1103099701a.pdf}
\IMGanswer{1103099701b.pdf}
\IMGanswer{1103099701c.pdf}
\IMGanswer{1103099701d.pdf}\End

\Data\ID{1103099702}
\Updated{2021-04-12 08:04:14}
\Level{A}
\SubArea{Logarithmic Functions}
\LANGen{
\Question{
One of the given graphs is not the graph of a logarithmic function. Choose this graph.}
{}{}{}{}{}{}}
\LANGpl{
\Question{
Który z poniższych wykresów nie przedstawia funkcji logarytmicznej?}
{}{}{}{}{}{}}
\LANGcs{
\Question{
Který z následujících grafů není grafem logaritmické funkce?}
{}{}{}{}{}{}}
\LANGsk{
\Question{
Ktorý z nasledujúcich grafov nie je graf logaritmickej funkcie?}
{}{}{}{}{}{}}
\LANGes{
\Question{
De las gráficas dadas, elige la de una función logarítmica.}
{}{}{}{}{}{}}

\IMGanswer{1103099702a_0.pdf}
\IMGanswer{1103099702b_0.pdf}
\IMGanswer{1103099702c_0.pdf}
\IMGanswer{1103099702d_0.pdf}\End

\Data\ID{1103099801}
\Updated{2020-11-12 01:11:41}
\Level{A}
\SubArea{Logarithmic Functions}
\IMG{1103099801.pdf}
\LANGen{
\Question{
The function \( f(x)=\log_2(x+3) \) is graphed below. Find the vertical asymptote of \( f \).}
{\( x=-3 \)}{\( x=3 \)}{\( y=-3 \)}{\( y=3 \)}{}{}}
\LANGpl{
\Question{
Funkcję  \( f(x)=\log_2(x+3) \) przedstawiono poniżej. Wyznacz pionową asymptotę \( f \).}
{\( x=-3 \)}{\( x=3 \)}{\( y=-3 \)}{\( y=3 \)}{}{}}
\LANGcs{
\Question{
Na obrázku je graf funkce \( f(x)=\log_2(x+3) \). Která rovnice je rovnicí svislé asymptoty dané funkce \( f \)?}
{\( x=-3 \)}{\( x=3 \)}{\( y=-3 \)}{\( y=3 \)}{}{}}
\LANGsk{
\Question{
Na obrázku je graf funkcie \( f(x)=\log_2(x+3) \). Ktorá rovnica je rovnicou zvislej asymptoty danej funkcie \( f \)?}
{\( x=-3 \)}{\( x=3 \)}{\( y=-3 \)}{\( y=3 \)}{}{}}
\LANGes{
\Question{
La gráfica de abajo representa la función \( f(x)=\log_2(x+3) \). Halla la asíntota vertical de la función \( f \).}
{\( x=-3 \)}{\( x=3 \)}{\( y=-3 \)}{\( y=3 \)}{}{}}
\End

\Data\ID{1103099802}
\Updated{2020-07-15 03:07:12}
\Level{A}
\SubArea{Logarithmic Functions}
\IMG{1103099802.pdf}
\LANGen{
\Question{
The function \( f(x)=\log_{\frac12} (x+3) \) is graphed below. Which of the following equations is the equation of the vertical asymptote of \( f \)?}
{\( x = -3  \)}{\( x = 3  \)}{\( y = -3  \)}{\( y = 3  \)}{}{}}
\LANGpl{
\Question{
Funkcję \( f(x)=\log_{\frac12} (x+3) \) przedstawiono poniżej. Które z poniższych równań jest równaniem pionowej asymptoty \( f \)?}
{\( x = -3  \)}{\( x = 3  \)}{\( y = -3  \)}{\( y = 3  \)}{}{}}
\LANGcs{
\Question{
Na obrázku je graf funkce \( f(x)=\log_{\frac12} (x+3) \). Která rovnice je rovnicí svislé asymptoty dané funkce \( f \)?}
{\( x = -3  \)}{\( x = 3  \)}{\( y = -3  \)}{\( y = 3  \)}{}{}}
\LANGsk{
\Question{
Na obrázku je graf funkcie \( f(x)=\log_{\frac12} (x+3) \). Ktorá rovnica je rovnicou zvislej asymptoty danej funkcie \( f \)?}
{\( x = -3  \)}{\( x = 3  \)}{\( y = -3  \)}{\( y = 3  \)}{}{}}
\LANGes{
\Question{
La gráfica representa la función \( f(x)=\log_{\frac12} (x+3) \). ¿Cuál de las siguientes ecuaciones es la ecuación de la asíntota vertical de la función \( f \)?}
{\( x = -3  \)}{\( x = 3  \)}{\( y = -3  \)}{\( y = 3  \)}{}{}}
\End

\Data\ID{1103099805}
\Updated{2021-04-12 09:04:08}
\Level{A}
\SubArea{Logarithmic Functions}
\IMG{1103099805.pdf}
\LANGen{
\Question{
Below, we are given the graph of the logarithmic function \( f(x)=\log_3 x \).  Use the graph of \( f \) as help to identify which of the next four given graphs is the graph of \( g(x)=\log_3 x-2 \).}
{}{}{}{}{}{}}
\LANGpl{
\Question{
Poniżej przedstawiono wykres funkcji logarytmicznej \( f(x)=\log_3 x \).  Wykorzystując wykres \( f \) określ, który z kolejnych czterech wykresów przedstawia funkcję \( g(x)=\log_3 x-2 \).}
{}{}{}{}{}{}}
\LANGcs{
\Question{
Na obrázku je graf logaritmické funkce \( f(x)=\log_3 x \).  Pomocí grafu funkce \( f \) identifikujte, který z následujících čtyř grafů je grafem funkce \( g(x)=\log_3 x-2 \).}
{}{}{}{}{}{}}
\LANGsk{
\Question{
Na obrázku je graf logaritmickej funkcie \( f(x)=\log_3 x \).  Pomocou grafu funkcie \( f \) identifikujte, ktorý z nasledujúcich štyroch grafov je graf funkcie \( g(x)=\log_3 x-2 \).}
{}{}{}{}{}{}}
\LANGes{
\Question{
Abajo aparece la gráfica de la función logarítmica \( f(x)=\log_3 x \).  Usando esta gráfica de la función \( f \), identifica cuál de las gráficas es la gráfica de la función \( g(x)=\log_3 x-2 \).}
{}{}{}{}{}{}}

\IMGanswer{1103099805a.pdf}
\IMGanswer{1103099805b.pdf}
\IMGanswer{1103099805c.pdf}
\IMGanswer{1103099805d.pdf}\End

\Data\ID{1103099806}
\Updated{2021-04-12 09:04:38}
\Level{A}
\SubArea{Logarithmic Functions}
\IMG{1103099806.pdf}
\LANGen{
\Question{
Below,  we are given the graph of the logarithmic function \( f(x)=\log_7x\).  Use the graph of \( f \) as help to identify which of the next four given graphs is the graph of \( g(x)=\log_7 x+2 \).}
{}{}{}{}{}{}}
\LANGpl{
\Question{
Poniżej przedstawiono wykres funkcji logarytmicznej  \( f(x)=\log_7x\).  Wykorzystując wykres \( f \) określ, który z czterech kolejnych wykresów obrazuje funkcję \( g(x)=\log_7 x+2 \).}
{}{}{}{}{}{}}
\LANGcs{
\Question{
Na obrázku je graf logaritmické funkce \( f(x)=\log_7x\). Pomocí grafu funkce \( f \) identifikujte, který z následujících čtyř grafů je grafem funkce \( g(x)=\log_7 x+2 \).}
{}{}{}{}{}{}}
\LANGsk{
\Question{
Na obrázku je graf logaritmickej funkcie \( f(x)=\log_7x\). Pomocou grafu funkcie \( f \) identifikujte, ktorý z nasledujúcich štyroch grafov je graf funkcie \( g(x)=\log_7 x+2 \).}
{}{}{}{}{}{}}
\LANGes{
\Question{
En la imagen podemos ver la gráfica de la función logarítmica \( f(x)=\log_7x\).  Usando esta gráfica de la función \( f \), identifica cuál de las gráficas es la gráfica de la función \( g(x)=\log_7 x+2 \).}
{}{}{}{}{}{}}

\IMGanswer{1103099806a.pdf}
\IMGanswer{1103099806b.pdf}
\IMGanswer{1103099806c.pdf}
\IMGanswer{1103099806d.pdf}\End

\Data\ID{1103099807}
\Updated{2021-04-12 09:04:12}
\Level{A}
\SubArea{Logarithmic Functions}
\IMG{1103099807.pdf}
\LANGen{
\Question{
Below, we are given the graph of the logarithmic function \( f(x)=\log_{\frac13} x \).  Use the graph of \( f \) as help to identify which of the next four given graphs is the graph of \( g(x)=\log_{\frac13}(x-2) \).}
{}{}{}{}{}{}}
\LANGpl{
\Question{
Poniżej przedstawiono wykres funkcji logarytmicznej \( f(x)=\log_{\frac13} x \).  Wykorzystując  \( f \) określ, który z czterech kolejnych wykresów obrazuje funkcję \( g(x)=\log_{\frac13}(x-2) \).}
{}{}{}{}{}{}}
\LANGcs{
\Question{
Na obrázku je graf logaritmické funkce \( f(x)=\log_{\frac13} x \). Pomocí grafu funkce \( f \) identifikujte, který z následujících čtyř grafů je grafem funkce \( g(x)=\log_{\frac13}(x-2) \).}
{}{}{}{}{}{}}
\LANGsk{
\Question{
Na obrázku je graf logaritmickej funkcie \( f(x)=\log_{\frac13} x \). Pomocou grafu funkcie \( f \) identifikujte, ktorý z nasledujúcich štyroch grafov je graf funkcie \( g(x)=\log_{\frac13}(x-2) \).}
{}{}{}{}{}{}}
\LANGes{
\Question{
En la imagen podemos ver la gráfica de la función logarítmica \( f(x)=\log_{\frac13} x \).  Usando la gráfica de la función \( f \), identifica cuál de las gráficas es la gráfica de la función \( g(x)=\log_{\frac13}(x-2) \).}
{}{}{}{}{}{}}

\IMGanswer{1103099807a.pdf}
\IMGanswer{1103099807b.pdf}
\IMGanswer{1103099807c.pdf}
\IMGanswer{1103099807d.pdf}\End

\Data\ID{1103099808}
\Updated{2021-04-12 09:04:25}
\Level{A}
\SubArea{Logarithmic Functions}
\IMG{1103099808_0.pdf}
\LANGen{
\Question{
Below,  we are given the graph of the logarithmic function \( f(x)=\log (x) \).  Use the graph of \( f \) as help to identify which of the next four given graphs is the graph of  \( g(x)=-4\log(x) \).}
{}{}{}{}{}{}}
\LANGpl{
\Question{
Poniżej przedstawiono wykres funkcji logarytmicznej \( f(x)=\log (x) \).  Wykorzystując wykres \( f \) określ, który z czterech kolejnych wykresów obrazuje funkcję   \( g(x)=-4\log(x) \).}
{}{}{}{}{}{}}
\LANGcs{
\Question{
Na obrázku je graf logaritmické funkce \( f(x)=\log (x) \). Pomocí grafu funkce \( f \) identifikujte, který z následujících čtyř grafů je grafem funkce \( g(x)=-4\log(x) \).}
{}{}{}{}{}{}}
\LANGsk{
\Question{
Na obrázku je graf logaritmickej funkcie \( f(x)=\log (x) \). Pomocou grafu funkcie \( f \) identifikujte, ktorý z nasledujúcich štyroch grafov je graf funkcie \( g(x)=-4\log(x) \).}
{}{}{}{}{}{}}
\LANGes{
\Question{
En la imagen podemos ver la gráfica de la función logarítmica \( f(x)=\log (x) \).  Usando esta gráfica de la función \( f \), identifica cuál de las cuatro gráficas es la gráfica de la función \( g(x)=-4\log(x) \).}
{}{}{}{}{}{}}

\IMGanswer{1103099808a_0.pdf}
\IMGanswer{1103099808b_0.pdf}
\IMGanswer{1103099808c_0.pdf}
\IMGanswer{1103099808d_0.pdf}\End

\Data\ID{1103099809}
\Updated{2020-07-15 03:07:10}
\Level{A}
\SubArea{Logarithmic Functions}
\IMG{1103099809.pdf}
\LANGen{
\Question{
Below,  we are given the graph of the logarithmic function \( f(x)=\log_4(x-3) \).  Use the graph of \( f \) as help to identify which of the next four given graphs is the graph of  \( g(x)=\log_4(3-x) \).}
{}{}{}{}{}{}}
\LANGpl{
\Question{
Poniżej przedstawiono wykres funkcji logarytmicznej \( f(x)=\log_4(x-3) \).  Wykorzystując wykres \( f \) określ, który z czterech kolejnych wykresów obrazuje funkcję \( g(x)=\log_4(3-x) \).}
{}{}{}{}{}{}}
\LANGcs{
\Question{
Na obrázku je graf logaritmické funkce \( f(x)=\log_4(x-3) \). Pomocí grafu funkce \( f \) identifikujte, který z následujících čtyř grafů je grafem funkce \( g(x)=\log_4(3-x) \).}
{}{}{}{}{}{}}
\LANGsk{
\Question{
Na obrázku je graf logaritmickej funkcie \( f(x)=\log_4(x-3) \). Pomocou grafu funkcie \( f \) identifikujte, ktorý z nasledujúcich štyroch grafov je graf funkcie \( g(x)=\log_4(3-x) \).}
{}{}{}{}{}{}}
\LANGes{
\Question{
En la imagen podemos ver la gráfica de la función logarítmica \( f(x)=\log_4(x-3) \).  Usando la gráfica de la función \( f \) identifica cuál de las siguientes cuatro gráficas es la gráfica de la función \( g(x)=\log_4(3-x) \).}
{}{}{}{}{}{}}

\IMGanswer{1103099809a.pdf}
\IMGanswer{1103099809b.pdf}
\IMGanswer{1103099809c.pdf}
\IMGanswer{1103099809d.pdf}\End

\Data\ID{1103099810}
\Updated{2020-11-12 01:11:28}
\Level{A}
\SubArea{Logarithmic Functions}
\LANGen{
\Question{
Which of the following graphs is the graph of the function \( f(x)=2-\log_{0.4}x \)?}
{}{}{}{}{}{}}
\LANGpl{
\Question{
Który z poniższych wykresów jest wykresem funkcji \( f(x)=2-\log_{0{,}4}x \)?}
{}{}{}{}{}{}}
\LANGcs{
\Question{
Který z následujících grafů je grafem funkce \( f(x)=2-\log_{0{,}4}x \)?}
{}{}{}{}{}{}}
\LANGsk{
\Question{
Ktorý z nasledujúcich grafov je graf funkcie \( f(x)=2-\log_{0{,}4}x \)?}
{}{}{}{}{}{}}
\LANGes{
\Question{
¿Cuál de las siguientes gráficas es la gráfica de la función \( f(x)=2-\log_{0.4}x \)?}
{}{}{}{}{}{}}

\IMGanswer{1103099810a_0.pdf}
\IMGanswer{1103099810b_0.pdf}
\IMGanswer{1103099810c_0.pdf}
\IMGanswer{1103099810d_0.pdf}\End

\Data\ID{1103099811}
\Updated{2020-11-12 01:11:35}
\Level{A}
\SubArea{Logarithmic Functions}
\LANGen{
\Question{
Which of the following graphs is the graph of the function \( f(x)=2\log_2x \)?}
{}{}{}{}{}{}}
\LANGpl{
\Question{
Który z podanych wykresów jest wykresem funkcji \( f(x)=2\log_2x \)?}
{}{}{}{}{}{}}
\LANGcs{
\Question{
Který z následujících grafů je grafem funkce \( f(x)=2\log_2x \)?}
{}{}{}{}{}{}}
\LANGsk{
\Question{
Ktorý z nasledujúcich grafov je graf funkcie \( f(x)=2\log_2x \)?}
{}{}{}{}{}{}}
\LANGes{
\Question{
¿Cuál de las siguientes gráficas es la gráfica de la función \( f(x)=2\log_2x \)?}
{}{}{}{}{}{}}

\IMGanswer{1103099811a_0.pdf}
\IMGanswer{1103099811b_0.pdf}
\IMGanswer{1103099811c_0.pdf}
\IMGanswer{1103099811d_0.pdf}\End

\Data\ID{2000000601}
\Updated{2021-12-01 03:12:14}
\Level{A}
\SubArea{Logarithmic Functions}
\LANGen{
\Question{
The graph of  \(f\colon y=\log_{2}x+2\) is obtained from the graph of \(g\colon y=\log_{2}x\) by shifting the graph of \(g\) by:}
{\(2\) units up.}{\(2\) units down.}{\(2\) units right.}{\(2\) units left.}{}{}}
\LANGpl{
\Question{
Wykres funkcji  \(f\colon y=\log_{2}x+2\) powstaje po przesunięciu wykresu funkcji \(g\colon y=\log_{2}x\):}
{o \(2\) jednostki w górę.}{o \(2\) jednostki w dół.}{o \(2\) jednostki w prawo.}{o \(2\) jednostki w lewo.}{}{}}
\LANGcs{
\Question{
Graf \(f\colon y=\log_{2}x+2\) získáme z grafu \(g\colon y=\log_{2}x\) posunutím grafu \(g\) o:}
{\(2\) jednotky nahoru.}{\(2\) jednotky dolů.}{\(2\) jednotky doprava.}{\(2\) jednotky doleva.}{}{}}
\LANGsk{
\Question{
Graf \(f\colon y=\log_{2}x+2\) získame z grafu \(g\colon y=\log_{2}x\) posunutím grafu \(g\) o:}
{\(2\) jednotky hore.}{\(2\) jednotky dole.}{\(2\) jednotky vpravo.}{\(2\) jednotky vľavo.}{}{}}
\LANGes{
\Question{
¿Cómo podemos conseguir la gráfica de \(f\colon y=\log_{2}x+2\) utilizando la gráfica de \(g\colon y=\log_{2}x\)?}
{subiendo la gráfica de \(g\) \(2\) unidades}{bajando la gráfica de \(g\) \(2\) unidades}{moviendo la gráfica de \(g\) \(2\) unidades a la derecha}{moviendo la gráfica de \(g\) \(2\) unidades a la izquierda}{}{}}
\End

\Data\ID{2000000602}
\Updated{2021-12-01 03:12:22}
\Level{A}
\SubArea{Logarithmic Functions}
\LANGen{
\Question{
The graph of which of the following functions contains the point \( \left[ \frac{1}{4} ; -1 \right]\)?}
{\(f:y=\log_{4}x\)}{\(f:y=\log_{\frac{1}{2}}x\)}{\(f:y=\log_{2}x\)}{\(f:y=\log_{\frac{1}{4}}x\)}{}{}}
\LANGpl{
\Question{
Punkt o współrzędnych \( \left[ \frac{1}{4} ; -1 \right]\) należy do wykresu funkcji opisanej wzorem:}
{\(f:y=\log_{4}x\)}{\(f:y=\log_{\frac{1}{2}}x\)}{\(f:y=\log_{2}x\)}{\(f:y=\log_{\frac{1}{4}}x\)}{}{}}
\LANGcs{
\Question{
Vyberte funkci, jejímuž grafu náleží bod \( \left[ \frac{1}{4} ; -1 \right]\).}
{\(f:y=\log_{4}x\)}{\(f:y=\log_{\frac{1}{2}}x\)}{\(f:y=\log_{2}x\)}{\(f:y=\log_{\frac{1}{4}}x\)}{}{}}
\LANGsk{
\Question{
Grafu ktorej z daných funkcií patrí bod \( \left[ \frac{1}{4} ; -1 \right]\)?}
{\(f:y=\log_{4}x\)}{\(f:y=\log_{\frac{1}{2}}x\)}{\(f:y=\log_{2}x\)}{\(f:y=\log_{\frac{1}{4}}x\)}{}{}}
\LANGes{
\Question{
Elige la función que contiene el punto \( \left[ \frac{1}{4} ; -1 \right]\).}
{\(f:y=\log_{4}x\)}{\(f:y=\log_{\frac{1}{2}}x\)}{\(f:y=\log_{2}x\)}{\(f:y=\log_{\frac{1}{4}}x\)}{}{}}
\End

\Data\ID{2000000603}
\Updated{2021-12-01 03:12:28}
\Level{A}
\SubArea{Logarithmic Functions}
\LANGen{
\Question{
Which of the following points does not lie on the graph of the function \(f: y=\log_{3}x\)?}
{\( [3;-1]\)}{\( [3;1]\)}{\( \left[ \frac{1}{3} ; -1 \right]\)}{\( [1;0]\)}{}{}}
\LANGpl{
\Question{
Który z punktów  nie należy do wykresu funkcji \(f: y=\log_{3}x\)?}
{\( [3;-1]\)}{\( [3;1]\)}{\( \left[ \frac{1}{3} ; -1 \right]\)}{\( [1;0]\)}{}{}}
\LANGcs{
\Question{
Který z následujících bodů nenáleží grafu funkce \(f: y=\log_{3}x\)?}
{\( [3;-1]\)}{\( [3;1]\)}{\( \left[ \frac{1}{3} ; -1 \right]\)}{\( [1;0]\)}{}{}}
\LANGsk{
\Question{
Ktorý z daných bodov neleži na grafe funkcie \(f: y=\log_{3}x\)?}
{\( [3;-1]\)}{\( [3;1]\)}{\( \left[ \frac{1}{3} ; -1 \right]\)}{\( [1;0]\)}{}{}}
\LANGes{
\Question{
Elige cuál de los puntos no está en la gráfica de la función \(f: y=\log_{3}x\).}
{\( [3;-1]\)}{\( [3;1]\)}{\( \left[ \frac{1}{3} ; -1 \right]\)}{\( [1;0]\)}{}{}}
\End

\Data\ID{2000000604}
\Updated{2021-12-01 03:12:33}
\Level{A}
\SubArea{Logarithmic Functions}
\LANGen{
\Question{
Which number does not belong to the domain of the function \(f: y=\log(x^2-4)\)?}
{\(x=\sqrt{3}\)}{\(x=-\sqrt{5}\)}{\(x=4\)}{\(x=-\sqrt{6}\)}{}{}}
\LANGpl{
\Question{
Która liczba nie należy do dziedziny funkcji: \(f: y=\log(x^2-4)\)?}
{\(x=\sqrt{3}\)}{\(x=-\sqrt{5}\)}{\(x=4\)}{\(x=-\sqrt{6}\)}{}{}}
\LANGcs{
\Question{
Které číslo nenáleží definičnímu oboru funkce \(f: y=\log(x^2-4)\)?}
{\(x=\sqrt{3}\)}{\(x=-\sqrt{5}\)}{\(x=4\)}{\(x=-\sqrt{6}\)}{}{}}
\LANGsk{
\Question{
Ktoré z daných čísel nepatrí do definičného oboru funkcie \(f: y=\log(x^2-4)\)?}
{\(x=\sqrt{3}\)}{\(x=-\sqrt{5}\)}{\(x=4\)}{\(x=-\sqrt{6}\)}{}{}}
\LANGes{
\Question{
Identifica el número que no pertenece al dominio de la función \(f: y=\log(x^2-4)\).}
{\(x=\sqrt{3}\)}{\(x=-\sqrt{5}\)}{\(x=4\)}{\(x=-\sqrt{6}\)}{}{}}
\End

\Data\ID{2000000605}
\Updated{2021-12-01 03:12:38}
\Level{A}
\SubArea{Logarithmic Functions}
\LANGen{
\Question{
Which of the following sets is the domain of the function \(f: y=\log(x^2+9)\)?}
{\( \mathbb{R}\)}{\( (-\infty;-3) \cup (3;+\infty)\)}{\( (-3;3)\)}{\( \mathbb{R} \setminus \{-3;3\}\)}{}{}}
\LANGpl{
\Question{
Który z poniższych zbiorów jest dziedziną funkcji \(f: y=\log(x^2+9)\)?}
{\( \mathbb{R}\)}{\( (-\infty;-3) \cup (3;+\infty)\)}{\( (-3;3)\)}{\( \mathbb{R} \setminus \{-3;3\}\)}{}{}}
\LANGcs{
\Question{
Která z následujících množin je definičním oborem funkce \(f: y=\log(x^2+9)\)?}
{\( \mathbb{R}\)}{\( (-\infty;-3) \cup (3;+\infty)\)}{\( (-3;3)\)}{\( \mathbb{R} \setminus \{-3;3\}\)}{}{}}
\LANGsk{
\Question{
Ktorá z nasledujúcich množin je definičným oborom funkcie \(f: y=\log(x^2+9)\)?}
{\( \mathbb{R}\)}{\( (-\infty;-3) \cup (3;+\infty)\)}{\( (-3;3)\)}{\( \mathbb{R} \setminus \{-3;3\}\)}{}{}}
\LANGes{
\Question{
Identifica el dominio de la función \(f: y=\log(x^2+9)\).}
{\( \mathbb{R}\)}{\( (-\infty;-3) \cup (3;+\infty)\)}{\( (-3;3)\)}{\( \mathbb{R} \setminus \{-3;3\}\)}{}{}}
\End

\Data\ID{2000000607}
\Updated{2021-12-05 12:12:41}
\Level{A}
\SubArea{Logarithmic Functions}
\LANGen{
\Question{
In how many points do the graphs of  \(f:y=-x+1\) and \(g:y=\log_{2}x\) intersect?}
{\(1\)}{\(0\)}{\(2\)}{\(3\)}{}{}}
\LANGpl{
\Question{
Ile punktów wspólnych mają wykresy funkcji \(f:y=-x+1\)  i  \(g:y=\log_{2}x\)?}
{\(1\)}{\(0\)}{\(2\)}{\(3\)}{}{}}
\LANGcs{
\Question{
V kolika bodech se protnou grafy \(f:y=-x+1\) a \(g:y=\log_{2}x\)?}
{\(1\)}{\(0\)}{\(2\)}{\(3\)}{}{}}
\LANGsk{
\Question{
V koľkých  bodoch sa pretnú grafy funkcií  \(f:y=-x+1\) a \(g:y=\log_{2}x\)?}
{\(1\)}{\(0\)}{\(2\)}{\(3\)}{}{}}
\LANGes{
\Question{
Identifica en cuántos puntos  intersecan las gráficas \(f:y=-x+1\) y \(g:y=\log_{2}x\).}
{\(1\)}{\(0\)}{\(2\)}{\(3\)}{}{}}
\End

\Data\ID{2010011007}
\Updated{2022-11-02 09:11:58}
\Level{A}
\SubArea{Logarithmic Functions}
\LANGen{
\Question{
In the following list identify the point that is not a point on the graph of the function.
\[f(x)=  -2\log _{2}x+3\]}
{\(\left[\frac12;1\right]\)}{\([2;1]\)}{\([4;-1]\)}{\([1;3]\)}{\(\left[\frac18;9\right]\)}{\(\left[\frac14;7\right]\)}}
\LANGpl{
\Question{
Który z punktów nie należy do wykresu funkcji?
\[f(x)=  -2\log _{2}x+3\]}
{\(\left[\frac12;1\right]\)}{\([2;1]\)}{\([4;-1]\)}{\([1;3]\)}{\(\left[\frac18;9\right]\)}{\(\left[\frac14;7\right]\)}}
\LANGcs{
\Question{
Který z bodů není bodem grafu funkce?
\[f(x)=  -2\log _{2}x+3\]}
{\(\left[\frac12;1\right]\)}{\([2;1]\)}{\([4;-1]\)}{\([1;3]\)}{\(\left[\frac18;9\right]\)}{\(\left[\frac14;7\right]\)}}
\LANGsk{
\Question{
Ktorým bodom neprechádza graf funkcie \(f(x)=  -2\log _{2}x+3\)?}
{\(\left[\frac12;1\right]\)}{\([2;1]\)}{\([4;-1]\)}{\([1;3]\)}{\(\left[\frac18;9\right]\)}{\(\left[\frac14;7\right]\)}}
\LANGes{
\Question{
En la siguiente lista, identifica el punto que no pertenece a la gráfica de la función.
\[f(x)=  -2\log _{2}x+3\]}
{\(\left[\frac12;1\right]\)}{\([2;1]\)}{\([4;-1]\)}{\([1;3]\)}{\(\left[\frac18;9\right]\)}{\(\left[\frac14;7\right]\)}}
\End

\Data\ID{2010011008}
\Updated{2022-11-02 09:11:58}
\Level{A}
\SubArea{Logarithmic Functions}
\LANGen{
\Question{
In the following list identify a positive expression.}
{\(\log _{0.5}3 -\log _{0.5}48\)}{\(\log _{0.5}16 +\log _{0.5}4\)}{\(\log _{3}9^3 -\log _{2}4^4\)}{\(\log _{5}\left(4^{-1}\right) +\log _{5}\frac4{125}\)}{}{}}
\LANGpl{
\Question{
Wskaż wyrażenie, którego wartość jest dodatnia.}
{\(\log _{0{,}5}3 -\log _{0{,}5}48\)}{\(\log _{0{,}5}16 +\log _{0{,}5}4\)}{\(\log _{3}9^3 -\log _{2}4^4\)}{\(\log _{5}\left(4^{-1}\right) +\log _{5}\frac4{125}\)}{}{}}
\LANGcs{
\Question{
Určete, který z výrazů je kladný.}
{\(\log _{0{,}5}3 -\log _{0{,}5}48\)}{\(\log _{0{,}5}16 +\log _{0{,}5}4\)}{\(\log _{3}9^3 -\log _{2}4^4\)}{\(\log _{5}\left(4^{-1}\right) +\log _{5}\frac4{125}\)}{}{}}
\LANGsk{
\Question{
Zistite, ktorý z daných výrazov je kladný.}
{\(\log _{0{,}5}3 -\log _{0{,}5}48\)}{\(\log _{0{,}5}16 +\log _{0{,}5}4\)}{\(\log _{3}9^3 -\log _{2}4^4\)}{\(\log _{5}\left(4^{-1}\right) +\log _{5}\frac4{125}\)}{}{}}
\LANGes{
\Question{
En la siguiente lista, identifica la expresión positiva.}
{\(\log _{0.5}3 -\log _{0.5}48\)}{\(\log _{0.5}16 +\log _{0.5}4\)}{\(\log _{3}9^3 -\log _{2}4^4\)}{\(\log _{5}\left(4^{-1}\right) +\log _{5}\frac4{125}\)}{}{}}
\End

\Data\ID{2010016009}
\Updated{2022-11-02 09:11:18}
\Level{A}
\SubArea{Logarithmic Functions}
\LANGen{
\Question{
Given the function \(f(x)=\log_2(x^2+4)\), evaluate \(f(2)\cdot f(0)\).}
{\(6\)}{\(32\)}{\(0\)}{\(24\)}{}{}}
\LANGpl{
\Question{
Mając daną funkcję \(f(x)=\log_2(x^2+4)\), oblicz \(f(2)\cdot f(0)\).}
{\(6\)}{\(32\)}{\(0\)}{\(24\)}{}{}}
\LANGcs{
\Question{
Je dána funkce \(f(x)=\log_2(x^2+4)\). Spočítejte součin \(f(2)\cdot f(0)\).}
{\(6\)}{\(32\)}{\(0\)}{\(24\)}{}{}}
\LANGsk{
\Question{
Dana je funkcia \(f(x)=\log_2(x^2+4)\), vypočítajte \(f(2)\cdot f(0)\).}
{\(6\)}{\(32\)}{\(0\)}{\(24\)}{}{}}
\LANGes{
\Question{
Dada la función \(f(x)=\log_2(x^2+4)\), evalúa \(f(2)\cdot f(0)\).}
{\(6\)}{\(32\)}{\(0\)}{\(24\)}{}{}}
\End

\Data\ID{2010016010}
\Updated{2022-11-02 09:11:18}
\Level{A}
\SubArea{Logarithmic Functions}
\LANGen{
\Question{
Given the function \(f(x)=\log_2(x+4)\), evaluate \(f(4)\cdot f(12)\).}
{\(12\)}{\(48\)}{\(128\)}{\(4\)}{}{}}
\LANGpl{
\Question{
Mając daną funkcję  \(f(x)=\log_2(x+4)\), oblicz \(f(4)\cdot f(12)\).}
{\(12\)}{\(48\)}{\(128\)}{\(4\)}{}{}}
\LANGcs{
\Question{
Je dána funkce \(f(x)=\log_2(x+4)\). Spočítejte součin \(f(4)\cdot f(12)\).}
{\(12\)}{\(48\)}{\(128\)}{\(4\)}{}{}}
\LANGsk{
\Question{
Daná je funkcia \(f(x)=\log_2(x+4)\), vypočítajte \(f(4)\cdot f(12)\).}
{\(12\)}{\(48\)}{\(128\)}{\(4\)}{}{}}
\LANGes{
\Question{
Dada la función \(f(x)=\log_2(x+4)\), evalúa \(f(4)\cdot f(12)\).}
{\(12\)}{\(48\)}{\(128\)}{\(4\)}{}{}}
\End

\Data\ID{9000003801}
\Updated{2020-07-15 03:07:34}
\Level{A}
\SubArea{Logarithmic Functions}
\IMG{000038_01-figure0.pdf}
\LANGen{
\Question{
Identify a possible analytic expression for the function graphed in the picture.}
{\(y =\log _{\frac{1} 
{2} }(x + 1) + 1\)}{\(y =\log _{\frac{1} 
{2} }(x - 1) + 1\)}{\(y =\log _{\frac{1} 
{2} }(x - 1) - 1\)}{\(y =\log _{\frac{1} 
{2} }(x + 1) - 1\)}{}{}}
\LANGpl{
\Question{
Określ wzór funkcji, której wykres przedstawiono poniżej.}
{\(y =\log _{\frac{1} 
{2} }(x + 1) + 1\)}{\(y =\log _{\frac{1} 
{2} }(x - 1) + 1\)}{\(y =\log _{\frac{1} 
{2} }(x - 1) - 1\)}{\(y =\log _{\frac{1} 
{2} }(x + 1) - 1\)}{}{}}
\LANGcs{
\Question{
Určete předpis funkce, jejíž graf je znázorněn na obrázku.}
{\(y =\log _{\frac{1} 
{2} }(x + 1) + 1\)}{\(y =\log _{\frac{1} 
{2} }(x - 1) + 1\)}{\(y =\log _{\frac{1} 
{2} }(x - 1) - 1\)}{\(y =\log _{\frac{1} 
{2} }(x + 1) - 1\)}{}{}}
\LANGsk{
\Question{
Určte predpis funkcie, ktorej graf je znázornený na obrázku.}
{\(y =\log _{\frac{1} 
{2} }(x + 1) + 1\)}{\(y =\log _{\frac{1} 
{2} }(x - 1) + 1\)}{\(y =\log _{\frac{1} 
{2} }(x - 1) - 1\)}{\(y =\log _{\frac{1} 
{2} }(x + 1) - 1\)}{}{}}
\LANGes{
\Question{
Identifica la posible expresión de la función representada en la imagen.}
{\(y =\log _{\frac{1} 
{2} }(x + 1) + 1\)}{\(y =\log _{\frac{1} 
{2} }(x - 1) + 1\)}{\(y =\log _{\frac{1} 
{2} }(x - 1) - 1\)}{\(y =\log _{\frac{1} 
{2} }(x + 1) - 1\)}{}{}}
\End

\Data\ID{9000003802}
\Updated{2023-04-03 02:04:02}
\Level{A}
\SubArea{Logarithmic Functions}
\LANGen{
\Question{
In the following list identify a function with graph through the points
\([5;0]\) and
\([-1;-2]\).}
{\(y =\log _{2}(x + 3) - 3\)}{\(y =\log _{5}(10 - x) - 1\)}{\(y =\log _{3}(4 + x) - 2\)}{\(y = 2 -\log _{3}(4 + x)\)}{\(y = 3 -\log _{2}(x + 3)\)}{\(y = 1 -\log _{5}(10 - x)\)}}
\LANGpl{
\Question{
Wybierz funkcję, której wykres przechodzi przez punkty
\([5;0]\) i
\([-1;-2]\).}
{\(y =\log _{2}(x + 3) - 3\)}{\(y =\log _{5}(10 - x) - 1\)}{\(y =\log _{3}(4 + x) - 2\)}{\(y = 2 -\log _{3}(4 + x)\)}{\(y = 3 -\log _{2}(x + 3)\)}{\(y = 1 -\log _{5}(10 - x)\)}}
\LANGcs{
\Question{
Vyberte předpis funkce, jejíž graf prochází body
\([5;0]\) a
\([-1;-2]\).}
{\(y =\log _{2}(x + 3) - 3\)}{\(y =\log _{5}(10 - x) - 1\)}{\(y =\log _{3}(4 + x) - 2\)}{\(y = 2 -\log _{3}(4 + x)\)}{\(y = 3 -\log _{2}(x + 3)\)}{\(y = 1 -\log _{5}(10 - x)\)}}
\LANGsk{
\Question{
Vyberte predpis funkcie, ktorej graf prechádza bodmi
\([5;0]\) a
\([-1;-2]\).}
{\(y =\log _{2}(x + 3) - 3\)}{\(y =\log _{5}(10 - x) - 1\)}{\(y =\log _{3}(4 + x) - 2\)}{\(y = 2 -\log _{3}(4 + x)\)}{\(y = 3 -\log _{2}(x + 3)\)}{\(y = 1 -\log _{5}(10 - x)\)}}
\LANGes{
\Question{
Identifica la función que pasa por los puntos 
\([5;0]\) y
\([-1;-2]\).}
{\(y =\log _{2}(x + 3) - 3\)}{\(y =\log _{5}(10 - x) - 1\)}{\(y =\log _{3}(4 + x) - 2\)}{\(y = 2 -\log _{3}(4 + x)\)}{\(y = 3 -\log _{2}(x + 3)\)}{\(y = 1 -\log _{5}(10 - x)\)}}
\End

\Data\ID{9000003804}
\Updated{2022-04-23 05:04:42}
\Level{A}
\SubArea{Logarithmic Functions}
\LANGen{
\Question{
In the following list identify the point that is not a point on the graph of the function:
\[f(x) = 1 -\log _{3}x\]}
{\([0;1]\)}{\([3;0]\)}{\(\left [\frac{1}
{9};3\right ]_{}\)}{\([1;1]\)}{\(\left [\frac{1}
{3};2\right ]\)}{\([9;-1]\)}}
\LANGpl{
\Question{
Który z poniższych punktów nie należy do wykresu funkcji 
\(f\colon y = 1 -\log _{3}x\)?}
{\([0;1]\)}{\([3;0]\)}{\(\left [\frac{1}
{9};3\right ]_{}\)}{\([1;1]\)}{\(\left [\frac{1}
{3};2\right ]\)}{\([9;-1]\)}}
\LANGcs{
\Question{
Kterým bodem neprochází graf funkce
\(f\colon y = 1 -\log _{3}x\)?}
{\([0;1]\)}{\([3;0]\)}{\(\left [\frac{1}
{9};3\right ]_{}\)}{\([1;1]\)}{\(\left [\frac{1}
{3};2\right ]\)}{\([9;-1]\)}}
\LANGsk{
\Question{
Ktorým bodom neprechádza graf funkcie
\(f\colon y = 1 -\log _{3}x\)?}
{\([0;1]\)}{\([3;0]\)}{\(\left [\frac{1}
{9};3\right ]_{}\)}{\([1;1]\)}{\(\left [\frac{1}
{3};2\right ]\)}{\([9;-1]\)}}
\LANGes{
\Question{
Identifica el punto por el que no pasa la gráfica de la función 
\(f\colon y = 1 -\log _{3}x\).}
{\([0;1]\)}{\([3;0]\)}{\(\left [\frac{1}
{9};3\right ]_{}\)}{\([1;1]\)}{\(\left [\frac{1}
{3};2\right ]\)}{\([9;-1]\)}}
\End

\Data\ID{9000003807}
\Updated{2022-04-23 05:04:42}
\Level{A}
\SubArea{Logarithmic Functions}
\LANGen{
\Question{
In the following list identify a negative expression.}
{\(\log _{0.1}20 -\log _{0.1}0.2\)}{\(\log _{3}9^{2.5} -\log _{4}4^{0.5}\)}{\(\log _{4}16^{\frac{3} 
{2} } +\log _{3}3^{\frac{1} 
{4} }\)}{\(\log _{3}7 +\log _{3}\frac{81} 
 {7} \)}{}{}}
\LANGpl{
\Question{
Które z poniższych wyrażeń jest ujemne?}
{\(\log _{0.1}20 -\log _{0.1}0.2\)}{\(\log _{3}9^{2.5} -\log _{4}4^{0.5}\)}{\(\log _{4}16^{\frac{3} 
{2} } +\log _{3}3^{\frac{1} 
{4} }\)}{\(\log _{3}7 +\log _{3}\frac{81} 
 {7} \)}{}{}}
\LANGcs{
\Question{
Určete, který výraz s logaritmy nabývá záporné hodnoty.}
{\(\log _{0{,}1}20 -\log _{0{,}1}0{,}2\)}{\(\log _{3}9^{2{,}5} -\log _{4}4^{0{,}5}\)}{\(\log _{4}16^{\frac{3} 
{2} } +\log _{3}3^{\frac{1} 
{4} }\)}{\(\log _{3}7 +\log _{3}\frac{81} 
 {7} \)}{}{}}
\LANGsk{
\Question{
Zistite, ktorý z daných výrazov je záporný.}
{\(\log _{0{,}1}20 -\log _{0{,}1}0{,}2\)}{\(\log _{3}9^{2{,}5} -\log _{4}4^{0{,}5}\)}{\(\log _{4}16^{\frac{3} 
{2} } +\log _{3}3^{\frac{1} 
{4} }\)}{\(\log _{3}7 +\log _{3}\frac{81} 
 {7} \)}{}{}}
\LANGes{
\Question{
Identifica una expresión negativa entre las siguentes:}
{\(\log _{0.1}20 -\log _{0.1}0.2\)}{\(\log _{3}9^{2.5} -\log _{4}4^{0.5}\)}{\(\log _{4}16^{\frac{3} 
{2} } +\log _{3}3^{\frac{1} 
{4} }\)}{\(\log _{3}7 +\log _{3}\frac{81} 
 {7} \)}{}{}}
\End

\Data\ID{9000004902}
\Updated{2023-04-03 02:04:25}
\Level{A}
\SubArea{Logarithmic Functions}
\LANGen{
\Question{
Find the domain of the function \(f\colon y =\log _{\frac{1} 
{3} }(9 - x^{2})\).}
{\(\mathrm{Dom}(f) = (-3;3)\)}{\(\mathrm{Dom}(f) =\mathbb{R}\setminus \{3\}\)}{\(\mathrm{Dom}(f) = (-\infty ;3)\)}{\(\mathrm{Dom}(f) = (3;\infty )\)}{\(\mathrm{Dom}(f) = (-\infty ;-3)\cup (3;\infty )\)}{}}
\LANGpl{
\Question{
Wyznacz dziedzinę funkcji  \(f\colon y =\log _{\frac{1} 
{3} }(9 - x^{2})\).}
{\(\mathrm{Dom}(f) = (-3;3)\)}{\(\mathrm{Dom}(f) =\mathbb{R}\setminus \{3\}\)}{\(\mathrm{Dom}(f) = (-\infty ;3)\)}{\(\mathrm{Dom}(f) = (3;\infty )\)}{\(\mathrm{Dom}(f) = (-\infty ;-3)\cup (3;\infty )\)}{}}
\LANGcs{
\Question{
Definiční obor funkce \(f\colon y =\log _{\frac{1} 
{3} }(9 - x^{2})\) 
je:}
{\(D(f) = (-3;3)\)}{\(D(f) =\mathbb{R}\setminus \{3\}\)}{\(D(f) = (-\infty ;3)\)}{\(D(f) = (3;\infty )\)}{\(D(f) = (-\infty ;-3)\cup (3;\infty )\)}{}}
\LANGsk{
\Question{
Definičným oborom funkcie \(f\colon y =\log _{\frac{1} 
{3} }(9 - x^{2})\) 
je:}
{\(D(f) = (-3;3)\)}{\(D(f) =\mathbb{R}\setminus \{3\}\)}{\(D(f) = (-\infty ;3)\)}{\(D(f) = (3;\infty )\)}{\(D(f) = (-\infty ;-3)\cup (3;\infty )\)}{}}
\LANGes{
\Question{
Halla el dominio de la función \(f\colon y =\log _{\frac{1} 
{3} }(9 - x^{2})\).}
{\(\mathrm{Dom}(f) = (-3;3)\)}{\(\mathrm{Dom}(f) =\mathbb{R}\setminus \{3\}\)}{\(\mathrm{Dom}(f) = (-\infty ;3)\)}{\(\mathrm{Dom}(f) = (3;\infty )\)}{\(\mathrm{Dom}(f) = (-\infty ;-3)\cup (3;\infty )\)}{}}
\End

\Data\ID{9000004903}
\Updated{2020-07-15 03:07:34}
\Level{A}
\SubArea{Logarithmic Functions}
\LANGen{
\Question{
Find the domain of the function \(f\colon y = \frac{3} 
{\log _{5}(x-4)}\).}
{\(\mathrm{Dom}(f) = (4;5)\cup (5;\infty )\)}{\(\mathrm{Dom}(f) = (0;\infty )\setminus \{4\}\)}{\(\mathrm{Dom}(f) = (-4;\infty )\setminus \{5\}\)}{\(\mathrm{Dom}(f) = (4;\infty )\)}{}{}}
\LANGpl{
\Question{
Wyznacz dziedzinę funkcji  \(f\colon y = \frac{3} 
{\log _{5}(x-4)}\).}
{\(\mathrm{Dom}(f) = (4;5)\cup (5;\infty )\)}{\(\mathrm{Dom}(f) = (0;\infty )\setminus \{4\}\)}{\(\mathrm{Dom}(f) = (-4;\infty )\setminus \{5\}\)}{\(\mathrm{Dom}(f) = (4;\infty )\)}{}{}}
\LANGcs{
\Question{
Definiční obor funkce \(f\colon y = \frac{3} 
{\log _{5}(x-4)}\) 
je:}
{\(D(f) = (4;5)\cup (5;\infty )\)}{\(D(f) = (0;\infty )\setminus \{4\}\)}{\(D(f) = (-4;\infty )\setminus \{5\}\)}{\(D(f) = (4;\infty )\)}{}{}}
\LANGsk{
\Question{
Definičným oborom funkcie \(f\colon y = \frac{3} 
{\log _{5}(x-4)}\) 
je:}
{\(D(f) = (4;5)\cup (5;\infty )\)}{\(D(f) = (0;\infty )\setminus \{4\}\)}{\(D(f) = (-4;\infty )\setminus \{5\}\)}{\(D(f) = (4;\infty )\)}{}{}}
\LANGes{
\Question{
Halla el dominio de la función \(f\colon y = \frac{3} 
{\log _{5}(x-4)}\).}
{\(\mathrm{Dom}(f) = (4;5)\cup (5;\infty )\)}{\(\mathrm{Dom}(f) = (0;\infty )\setminus \{4\}\)}{\(\mathrm{Dom}(f) = (-4;\infty )\setminus \{5\}\)}{\(\mathrm{Dom}(f) = (4;\infty )\)}{}{}}
\End

\Data\ID{9000004904}
\Updated{2020-07-15 03:07:34}
\Level{A}
\SubArea{Logarithmic Functions}
\LANGen{
\Question{
In the following list identify a function with a domain
\(\left (-\infty ; \frac{2} 
{3}\right ).\)}
{\(y =\log (2 - 3x)\)}{\(y =\log (3x - 2)\)}{\(y = -\log (3x - 2)\)}{\(y =\log (2x - 3)\)}{\(y =\log (3 - 2x)\)}{none of the above}}
\LANGpl{
\Question{
Dziedziną której z poniższych funkcji jest 
\(\left (-\infty ; \frac{2} 
{3}\right )\)?}
{\(y =\log (2 - 3x)\)}{\(y =\log (3x - 2)\)}{\(y = -\log (3x - 2)\)}{\(y =\log (2x - 3)\)}{\(y =\log (3 - 2x)\)}{żadne z powyższych}}
\LANGcs{
\Question{
Vyberte funkci, jejímž definičním oborem je interval
\(\left (-\infty ; \frac{2} 
{3}\right ).\)}
{\(y =\log (2 - 3x)\)}{\(y =\log (3x - 2)\)}{\(y = -\log (3x - 2)\)}{\(y =\log (2x - 3)\)}{\(y =\log (3 - 2x)\)}{žádná z uvedených funkcí}}
\LANGsk{
\Question{
Vyberte funkciu, ktorej definičný obor je interval
\(\left (-\infty ; \frac{2} 
{3}\right ).\)}
{\(y =\log (2 - 3x)\)}{\(y =\log (3x - 2)\)}{\(y = -\log (3x - 2)\)}{\(y =\log (2x - 3)\)}{\(y =\log (3 - 2x)\)}{žiadna z uvedených funkcií}}
\LANGes{
\Question{
Identifica la función cuyo dominio es el intervalo
\(\left (-\infty ; \frac{2} 
{3}\right ).\)}
{\(y =\log (2 - 3x)\)}{\(y =\log (3x - 2)\)}{\(y = -\log (3x - 2)\)}{\(y =\log (2x - 3)\)}{\(y =\log (3 - 2x)\)}{ninguna de las funciones dadas}}
\End

\Data\ID{9000004906}
\Updated{2020-07-15 03:07:34}
\Level{A}
\SubArea{Logarithmic Functions}
\IMG{000049_06-figure0.pdf}
\LANGen{
\Question{
Identify a possible analytic expression for the function
\(f\) 
graphed in the picture.}
{\(y =\log _{2}x\)}{\(y =\log _{0.2}x\)}{\(y =\log _{0.5}x\)}{\(y =\log _{5}x\)}{}{}}
\LANGpl{
\Question{
Wykres, której z poniższych funkcji
\(f\) 
został przedstawiony na rysunku poniżej?}
{\(y =\log _{2}x\)}{\(y =\log _{0.2}x\)}{\(y =\log _{0.5}x\)}{\(y =\log _{5}x\)}{}{}}
\LANGcs{
\Question{
Předpis funkce \(f\),
jejíž graf vidíte na obrázku, je:}
{\(y =\log _{2}x\)}{\(y =\log _{0{,}2}x\)}{\(y =\log _{0{,}5}x\)}{\(y =\log _{5}x\)}{}{}}
\LANGsk{
\Question{
Predpis funkcie \(f\),
ktorej graf vidíte na obrázku, je:}
{\(y =\log _{2}x\)}{\(y =\log _{0{,}2}x\)}{\(y =\log _{0{,}5}x\)}{\(y =\log _{5}x\)}{}{}}
\LANGes{
\Question{
Identifica una expresión analítica posible de la función \(f\) 
dada en la imagen.}
{\(y =\log _{2}x\)}{\(y =\log _{0.2}x\)}{\(y =\log _{0.5}x\)}{\(y =\log _{5}x\)}{}{}}
\End

\Data\ID{9000004909}
\Updated{2020-07-15 03:07:34}
\Level{A}
\SubArea{Logarithmic Functions}
\IMG{000049_09-figure0.pdf}
\LANGen{
\Question{
Identify a possible analytic expression for the function
\(g\) 
graphed in the picture.}
{\(y =\log _{3}(x + 2) - 1\)}{\(y =\log _{\frac{1} 
{3} }(x + 2) - 1\)}{\(y =\log _{3}(x - 2) + 1\)}{\(y =\log _{\frac{1} 
{3} }(x - 2) + 1\)}{}{}}
\LANGpl{
\Question{
Określ prawdopodobny wzór funkcji
\(g\), której wykres przedstawiono na rysunku poniżej.}
{\(y =\log _{3}(x + 2) - 1\)}{\(y =\log _{\frac{1} 
{3} }(x + 2) - 1\)}{\(y =\log _{3}(x - 2) + 1\)}{\(y =\log _{\frac{1} 
{3} }(x - 2) + 1\)}{}{}}
\LANGcs{
\Question{
Předpis funkce \(g\),
jejíž graf vidíte na obrázku, je:}
{\(y =\log _{3}(x + 2) - 1\)}{\(y =\log _{\frac{1} 
{3} }(x + 2) - 1\)}{\(y =\log _{3}(x - 2) + 1\)}{\(y =\log _{\frac{1} 
{3} }(x - 2) + 1\)}{}{}}
\LANGsk{
\Question{
Predpis funkcie \(g\),
ktorej graf vidíte na obrázku, je:}
{\(y =\log _{3}(x + 2) - 1\)}{\(y =\log _{\frac{1} 
{3} }(x + 2) - 1\)}{\(y =\log _{3}(x - 2) + 1\)}{\(y =\log _{\frac{1} 
{3} }(x - 2) + 1\)}{}{}}
\LANGes{
\Question{
Identifica la expresión analítica de la función 
\(g\) representada en la imagen.}
{\(y =\log _{3}(x + 2) - 1\)}{\(y =\log _{\frac{1} 
{3} }(x + 2) - 1\)}{\(y =\log _{3}(x - 2) + 1\)}{\(y =\log _{\frac{1} 
{3} }(x - 2) + 1\)}{}{}}
\End

\Data\ID{9100004907}
\Updated{2020-07-15 03:07:10}
\Level{A}
\SubArea{Logarithmic Functions}
\LANGen{
\Question{
Identify a possible graph of the function
\(f\colon y =\log _{0.2}x\).}
{}{}{}{}{}{}}
\LANGpl{
\Question{
Który z poniższych rysunków przedstawia wykres funkcji
\(f\colon y =\log _{0.2}x\)?}
{}{}{}{}{}{}}
\LANGcs{
\Question{
Na kterém z následujících obrázků je graf funkce
\(f\colon y =\log _{0{,}2}x\)?}
{}{}{}{}{}{}}
\LANGsk{
\Question{
Ktorý z následujúcich grafov je grafom funkcie 
\(f\colon y =\log _{0{,}2}x\)?}
{}{}{}{}{}{}}
\LANGes{
\Question{
Identifica la gráfica de la función 
\(f\colon y =\log _{0.2}x\).}
{}{}{}{}{}{}}

\IMGanswer{000049_07-figure2.pdf}
\IMGanswer{000049_07-figure0.pdf}
\IMGanswer{000049_07-figure1.pdf}
\IMGanswer{000049_07-figure3.pdf}\End

\Data\ID{1003118501}
\Updated{2020-07-15 03:07:12}
\Level{B}
\SubArea{Logarithmic Functions}
\LANGen{
\Question{
Given a  decreasing logarithmic function \( f(x)=\log_ax \), choose the correct statement about the base \( a \).}
{\( 0 < a < 1 \)}{\( a > 1 \)}{\( a=1 \)}{\( a < 1 \)}{}{}}
\LANGpl{
\Question{
Dana jest malejąca funkcja logarytmiczna \( f(x)=\log_ax \). Wybierz poprawne stwierdzenie dotyczące podstawy \( a \).}
{\( 0 < a < 1 \)}{\( a > 1 \)}{\( a=1 \)}{\( a < 1 \)}{}{}}
\LANGcs{
\Question{
Logaritmická funkce \( f(x)=\log_ax \) je klesající, když pro základ \( a \) platí:}
{\( 0 < a < 1 \)}{\( a > 1 \)}{\( a=1 \)}{\( a < 1 \)}{}{}}
\LANGsk{
\Question{
Logaritmická funkcia \( f(x)=\log_ax \) je klesajúca, ak pre základ \( a \) platí:}
{\( 0 < a < 1 \)}{\( a > 1 \)}{\( a=1 \)}{\( a < 1 \)}{}{}}
\LANGes{
\Question{
Dada la función logarítmica \( f(x)=\log_ax \) que es decreciente. Elige el enunciado correcto sobre la base \( a \).}
{\( 0 < a < 1 \)}{\( a > 1 \)}{\( a=1 \)}{\( a < 1 \)}{}{}}
\End

\Data\ID{1003118502}
\Updated{2021-04-12 08:04:11}
\Level{B}
\SubArea{Logarithmic Functions}
\LANGen{
\Question{
Given a  function \( f(x)=b\cdot \log_ax \), where \( a > 1 \) and \( b < 0 \), find the correct statement.}
{Function \( f \) is decreasing.}{Function \( f \) is increasing.}{Function \( f \) is bounded.}{Function \( f \) is bounded below.}{}{}}
\LANGpl{
\Question{
Dana jest funkcja \( f(x)=b\cdot \log_ax \), gdzie \( a > 1 \) and \( b < 0 \). Wybierz poprawne stwierdzenie.}
{Funkcja \( f \) jest malejąca.}{Funkcja \( f \) jest rosnąca.}{Funkcja \( f \) jest ograniczona.}{Funkcja \( f \) jest ograniczona z dołu.}{}{}}
\LANGcs{
\Question{
Je dána funkce \( f(x)=b\cdot \log_ax \), kde \( a > 1 \) a \( b < 0 \). Vyberte pravdivé tvrzení.}
{Funkce \( f \) je klesající.}{Funkce \( f \) je rostoucí.}{Funkce \( f \) je omezená.}{Funkce \( f \) je omezená zdola.}{}{}}
\LANGsk{
\Question{
Daná je funkcia \( f(x)=b\cdot \log_ax \), kde \( a > 1 \) a \( b < 0 \). Vyberte pravdivé tvrdenie.}
{Funkcia \( f \) je klesajúca.}{Funkcia \( f \) je rastúca.}{Funkcia \( f \) je ohraničená.}{Funkcia \( f \) je ohraničená zdola.}{}{}}
\LANGes{
\Question{
Dada la función \( f(x)=b\cdot \log_ax \), donde \( a > 1 \) y \( b < 0 \), halla el enunciado correcto.}
{La función \( f \) es decreciente.}{La función \( f \) es creciente.}{La función \( f \) está acotada.}{La función \( f \) está acotada inferiormente.}{}{}}
\End

\Data\ID{1003118503}
\Updated{2020-11-12 01:11:44}
\Level{B}
\SubArea{Logarithmic Functions}
\LANGen{
\Question{
Which of the following functions is increasing?}
{\( f\colon y=-\log_{\frac12}x \)}{\( f\colon y=-\log_2x \)}{\( f\colon y=-2\log_2x \)}{\( f\colon y=\log_{\frac12}x \)}{}{}}
\LANGpl{
\Question{
Która z poniższych funkcji jest rosnąca?}
{\( f\colon y=-\log_{\frac12}x \)}{\( f\colon y=-\log_2x \)}{\( f\colon y=-2\log_2x \)}{\( f\colon y=\log_{\frac12}x \)}{}{}}
\LANGcs{
\Question{
Která z následujících funkcí je rostoucí?}
{\( f\colon y=-\log_{\frac12}x \)}{\( f\colon y=-\log_2x \)}{\( f\colon y=-2\log_2x \)}{\( f\colon y=\log_{\frac12}x \)}{}{}}
\LANGsk{
\Question{
Ktorá z nasledujúcich funkcií je rastúca?}
{\( f\colon y=-\log_{\frac12}x \)}{\( f\colon y=-\log_2x \)}{\( f\colon y=-2\log_2x \)}{\( f\colon y=\log_{\frac12}x \)}{}{}}
\LANGes{
\Question{
¿Cuál de las siguientes funciones es creciente?}
{\( f\colon y=-\log_{\frac12}x \)}{\( f\colon y=-\log_2x \)}{\( f\colon y=-2\log_2x \)}{\( f\colon y=\log_{\frac12}x \)}{}{}}
\End

\Data\ID{1003136502}
\Updated{2021-04-12 08:04:46}
\Level{B}
\SubArea{Logarithmic Functions}
\LANGen{
\Question{
Consider the values 

\(\ \log_74;\) \(\log_{\frac47}{0.4};\) \(\log_40.7;\) \(\log_{\frac74}4;\) \(\log_{0.7}0.4;\) \(\log_{0.4}4;\) \(\log_7{0.7}.\ \)

Without using a calculator, determine how many of the given values are positive.}
{\( 4 \)}{\( 3 \)}{\( 2 \)}{\( 5 \)}{}{}}
\LANGpl{
\Question{
Rozważ wartości
\(\ \log_74;\) \(\log_{\frac47}{0{,}4};\) \(\log_40{,}7;\) \(\log_{\frac74}4;\) \(\log_{0{,}7}0{,}4;\) \(\log_{0{,}4}4;\) \(\log_7{0{,}7}.\ \)
Nie używając kalkulatora określ, ile spośród podanych wartości jest dodatnich.}
{\( 4 \)}{\( 3 \)}{\( 2 \)}{\( 5 \)}{}{}}
\LANGcs{
\Question{
Mějme hodnoty 

\(\ \log_74;\) \(\log_{\frac47}{0{,}4};\) \(\log_40{,}7;\) \(\log_{\frac74}4;\) \(\log_{0{,}7}0{,}4;\) \(\log_{0{,}4}4;\) \(\log_7{0{,}7}.\ \)

Bez použití kalkulačky určete, kolik z daných hodnot je kladných.}
{\( 4 \)}{\( 3 \)}{\( 2 \)}{\( 5 \)}{}{}}
\LANGsk{
\Question{
Bez použitia kalkulačky určte, koľko z nasledujúcich hodnôt je kladných.
\(\ \log_74;\) \(\log_{\frac47}{0{,}4};\) \(\log_40{,}7;\) \(\log_{\frac74}4;\) \(\log_{0{,}7}0{,}4;\) \(\log_{0{,}4}4;\) \(\log_7{0{,}7}.\ \)}
{\( 4 \)}{\( 3 \)}{\( 2 \)}{\( 5 \)}{}{}}
\LANGes{
\Question{
Teniendo los valores
\(\ \log_74;\) \(\log_{\frac47}{0.4};\) \(\log_40.7;\) \(\log_{\frac74}4;\) \(\log_{0.7}0.4;\) \(\log_{0.4}4;\) \(\log_7{0.7}.\ \)
Sin utilizar la calculadora, determina cuántos valores son positivos.}
{\( 4 \)}{\( 3 \)}{\( 2 \)}{\( 5 \)}{}{}}
\End

\Data\ID{1003136507}
\Updated{2021-04-12 08:04:17}
\Level{B}
\SubArea{Logarithmic Functions}
\LANGen{
\Question{
What are the values of the real parameter \( a \), that satisfy inequality \( \log_a\frac47 > \log_a \frac74 \)?}
{\( 0 < a < 1 \)}{\( a > 1 \)}{\( a < 1 \)}{\( a > 0 \)}{}{}}
\LANGpl{
\Question{
Jakie są wartości rzeczywistego parametru \( a \), który spełnia nierówność \( \log_a\frac47 > \log_a \frac74 \)?}
{\( 0 < a < 1 \)}{\( a > 1 \)}{\( a < 1 \)}{\( a > 0 \)}{}{}}
\LANGcs{
\Question{
Určete všechna reálná \( a \), pro něž platí \( \log_a\frac47 > \log_a \frac74 \).}
{\( 0 < a < 1 \)}{\( a > 1 \)}{\( a < 1 \)}{\( a > 0 \)}{}{}}
\LANGsk{
\Question{
Aké sú hodnoty reálneho parametra \( a \), ktoré vyhovujú nerovnici \( \log_a\frac47 > \log_a \frac74 \)?}
{\( 0 < a < 1 \)}{\( a > 1 \)}{\( a < 1 \)}{\( a > 0 \)}{}{}}
\LANGes{
\Question{
Halla los valores del parámetro real \( a \), para que sea cierta la siguiente desigualdad \( \log_a\frac47 > \log_a \frac74 \).}
{\( 0 < a < 1 \)}{\( a > 1 \)}{\( a < 1 \)}{\( a > 0 \)}{}{}}
\End

\Data\ID{1003136508}
\Updated{2021-04-12 08:04:56}
\Level{B}
\SubArea{Logarithmic Functions}
\LANGen{
\Question{
What are the values of the real parameter \( a \), that satisfy inequality \( \log_a7 < \log_a4 \)?}
{\( 0 < a < 1 \)}{\( a > 1 \)}{\( a < 1 \)}{\( a > 0 \)}{}{}}
\LANGpl{
\Question{
Jakie są wartości rzeczywistego parametru \( a \), który spełnia nierówność \( \log_a7 < \log_a4 \)?}
{\( 0 < a < 1 \)}{\( a > 1 \)}{\( a < 1 \)}{\( a > 0 \)}{}{}}
\LANGcs{
\Question{
Určete všechna reálná \( a \), pro něž platí \( \log_a7 < \log_a4 \)?}
{\( 0 < a < 1 \)}{\( a > 1 \)}{\( a < 1 \)}{\( a > 0 \)}{}{}}
\LANGsk{
\Question{
Aké sú hodnoty reálneho parametra \( a \), ktoré vyhovujú nerovnici \( \log_a7 < \log_a4 \)?}
{\( 0 < a < 1 \)}{\( a > 1 \)}{\( a < 1 \)}{\( a > 0 \)}{}{}}
\LANGes{
\Question{
Halla el valor del parámetro real \( a \) para el que la desigualdad \( \log_a7 < \log_a4 \) es cierta.}
{\( 0 < a < 1 \)}{\( a > 1 \)}{\( a < 1 \)}{\( a > 0 \)}{}{}}
\End

\Data\ID{1003136509}
\Updated{2020-07-15 03:07:21}
\Level{B}
\SubArea{Logarithmic Functions}
\LANGen{
\Question{
How many of the following inequalities is/are true?
\begin{align*}
\log_74 & > \log_711; & \log_{0.4} 0.7 &\leq \log_{0.4}3 \\ \log_{\frac17}4 &\geq \log_{\frac17}0.4; & \log_30.11 & < \log_36 \end{align*}}
{\( 1 \)}{\( 2 \)}{\( 3 \)}{\( 0 \)}{}{}}
\LANGpl{
\Question{
Ile z podanych nierówności jest prawdziwych? 
\begin{align*}
\log_74 & > \log_711; & \log_{0{,}4} 0{,}7 &\leq \log_{0{,}4}3 \\ \log_{\frac17}4 &\geq \log_{\frac17}0{,}4; & \log_30{,}11 & < \log_36 \end{align*}}
{\( 1 \)}{\( 2 \)}{\( 3 \)}{\( 0 \)}{}{}}
\LANGcs{
\Question{
Kolik z následujících nerovností je pravdivých?
\begin{align*}
\log_74 & > \log_711; & \log_{0{,}4} 0{,}7 &\leq \log_{0{,}4}3 \\ \log_{\frac17}4 &\geq \log_{\frac17}0{,}4; & \log_30{,}11 & < \log_36 \end{align*}}
{\( 1 \)}{\( 2 \)}{\( 3 \)}{\( 0 \)}{}{}}
\LANGsk{
\Question{
Koľko z nasledujúcich nerovností je pravdivých?
\begin{align*}
\log_74 & > \log_711; & \log_{0{,}4} 0{,}7 &\leq \log_{0{,}4}3 \\ \log_{\frac17}4 &\geq \log_{\frac17}0{,}4; & \log_30{,}11 & < \log_36 \end{align*}}
{\( 1 \)}{\( 2 \)}{\( 3 \)}{\( 0 \)}{}{}}
\LANGes{
\Question{
¿Cuántas de las siguientes desigualdades son verdaderas? 
\begin{align*}
\log_74 & > \log_711; & \log_{0.4} 0.7 &\leq \log_{0.4}3 \\ \log_{\frac17}4 &\geq \log_{\frac17}0.4; & \log_30.11 & < \log_36 \end{align*}}
{\( 1 \)}{\( 2 \)}{\( 3 \)}{\( 0 \)}{}{}}
\End

\Data\ID{1003136511}
\Updated{2020-07-15 03:07:21}
\Level{B}
\SubArea{Logarithmic Functions}
\LANGen{
\Question{
Identify which of the following relations is correct.}
{\( \log_28 > \log_2\frac72 > \log_21 > \log_2\frac47 > \log_2\frac17 \)}{\( \log_2\frac17 > \log_2\frac47 > \log_2 1 > \log_2\frac74 > \log_28 \)}{\( \log_28 > \log_2\frac74 > \log_2\frac47 > \log_2\frac17 > \log_21 \)}{\( \log_28 > \log_2\frac74 > \log_2\frac47 > \log_21 > \log_2\frac17 \)}{}{}}
\LANGpl{
\Question{
Określ, która z podanych zależności jest poprawna.}
{\( \log_28 > \log_2\frac72 > \log_21 > \log_2\frac47 > \log_2\frac17 \)}{\( \log_2\frac17 > \log_2\frac47 > \log_2 1 > \log_2\frac74 > \log_28 \)}{\( \log_28 > \log_2\frac74 > \log_2\frac47 > \log_2\frac17 > \log_21 \)}{\( \log_28 > \log_2\frac74 > \log_2\frac47 > \log_21 > \log_2\frac17 \)}{}{}}
\LANGcs{
\Question{
Určete, který z následujících vztahů je správný.}
{\( \log_28 > \log_2\frac72 > \log_21 > \log_2\frac47 > \log_2\frac17 \)}{\( \log_2\frac17 > \log_2\frac47 > \log_2 1 > \log_2\frac74 > \log_28 \)}{\( \log_28 > \log_2\frac74 > \log_2\frac47 > \log_2\frac17 > \log_21 \)}{\( \log_28 > \log_2\frac74 > \log_2\frac47 > \log_21 > \log_2\frac17 \)}{}{}}
\LANGsk{
\Question{
Určte, ktorý z nasledujúcich vzťahov je správny.}
{\( \log_28 > \log_2\frac72 > \log_21 > \log_2\frac47 > \log_2\frac17 \)}{\( \log_2\frac17 > \log_2\frac47 > \log_2 1 > \log_2\frac74 > \log_28 \)}{\( \log_28 > \log_2\frac74 > \log_2\frac47 > \log_2\frac17 > \log_21 \)}{\( \log_28 > \log_2\frac74 > \log_2\frac47 > \log_21 > \log_2\frac17 \)}{}{}}
\LANGes{
\Question{
Identifica cuál de las siguientes relaciones es correcta.}
{\( \log_28 > \log_2\frac72 > \log_21 > \log_2\frac47 > \log_2\frac17 \)}{\( \log_2\frac17 > \log_2\frac47 > \log_2 1 > \log_2\frac74 > \log_28 \)}{\( \log_28 > \log_2\frac74 > \log_2\frac47 > \log_2\frac17 > \log_21 \)}{\( \log_28 > \log_2\frac74 > \log_2\frac47 > \log_21 > \log_2\frac17 \)}{}{}}
\End

\Data\ID{1103118504}
\Updated{2021-02-22 04:02:31}
\Level{B}
\SubArea{Logarithmic Functions}
\IMG{1103118504.pdf}
\LANGen{
\Question{
Given the next  graph of  a logarithmic function \( f \), choose the correct list of properties of \( f \).}
{decreasing, asymptote at \( x=0 \), unbounded}{decreasing, asymptote at \( y=0 \), unbounded}{decreasing, asymptote at \( x=0 \), bounded below}{increasing, asymptote at \( y=0 \), bounded below}{}{}}
\LANGpl{
\Question{
Przedstawiony jest wykres funkcji logarytmicznej \( f \), wybierz listę właściwości tej funkcji \( f \).}
{malejąca, asymptota w \( x=0 \), nieograniczona}{malejąca, asymptota w \( y=0 \), nieograniczona}{malejąca, asymptota w \( x=0 \), ograniczona z dołu}{rosnąca, asymptota w \( y=0 \), ograniczona z dołu}{}{}}
\LANGcs{
\Question{
Která z následujících možností uvádí vlastnosti funkce \( f \), jejíž graf je na obrázku?}
{klesající, asymptota \( x=0 \), není omezená}{klesající, asymptota \( y=0 \), není omezená}{klesající, asymptota \( x=0 \), omezená zdola}{rostoucí, asymptota \( y=0 \), omezená zdola}{}{}}
\LANGsk{
\Question{
Daný je graf logaritmickej funkcie \( f \). Vyberte pravdivé tvrdenie o vlastnostiach funkcie \( f \).}
{klesajúca, asymptota \( x=0 \), neohraničená}{klesajúca, asymptota \( y=0 \), neohraničená}{klesajúca, asymptota \( x=0 \), ohraničená zdola}{rastúca, asymptota \( y=0 \), ohraničená zdola}{}{}}
\LANGes{
\Question{
Dada la función logarítmica \( f \), elige las propiedades verdaderas de la función \( f \).}
{decreciente, asíntota en \( x=0 \), no acotada}{decreciente, asíntota en \( y=0 \), no acotada}{decreciente, asíntota en \( x=0 \), acotada por abajo}{creciente, asíntota en \( y=0 \), acotada por abajo}{}{}}
\End

\Data\ID{1103118505}
\Updated{2021-04-12 09:04:10}
\Level{B}
\SubArea{Logarithmic Functions}
\LANGen{
\Question{
Suppose a logarithmic function \( f \) is unbounded, increasing and \( f \) has an asymptote at \( x=-2 \). Given properties of \( f \),  identify a possible graph of this function.}
{}{}{}{}{}{}}
\LANGpl{
\Question{
Załóżmy, że  funkcja logarytmiczna  \( f \) jest nieograniczona, rosnąca i \( f \) ma asymptotę w \( x=-2 \). Mając podane właściwości funkcji \( f \), wybierz możliwy wykres tej funkcji.}
{}{}{}{}{}{}}
\LANGcs{
\Question{
Který z následujících grafů je grafem logaritmické funkce \( f \), když předpokládáme, že funkce \( f \) není omezená, je rostoucí a má asymptotu \( x=-2 \)?}
{}{}{}{}{}{}}
\LANGsk{
\Question{
Ktorý z nasledujúcich grafov je graf logaritmickej funkcie \( f \), ak predpokladáme, že funkcia \( f \) je neohraničená, rastúca a má asymptotu \( x=-2 \)?}
{}{}{}{}{}{}}
\LANGes{
\Question{
Supongamos que la función logarítmica \( f \) no es acotada, es creciente y tiene una asíntota \( x=-2 \). Identifica cuál de las gráficas es la de esta función \( f \).}
{}{}{}{}{}{}}

\IMGanswer{1103118505a.pdf}
\IMGanswer{1103118505b.pdf}
\IMGanswer{1103118505c.pdf}
\IMGanswer{1103118505d.pdf}\End

\Data\ID{1103136501}
\Updated{2021-04-12 09:04:41}
\Level{B}
\SubArea{Logarithmic Functions}
\IMG{1103136501.pdf}
\LANGen{
\Question{
Consider the values
\( \log_{0.2}3;\) \(\log_20.3;\)  \(\log_{0.3}1;\) \(\log_23;\) \(\log_{\frac32}\frac23;\) \(\log_{\frac23}2;\) \(\log_{0.3}\frac32.\)
 Determine how many of the given values are negative. Use the graph of a logarithmic function given below.}
{\( 5 \)}{\( 4 \)}{\( 3 \)}{\( 2 \)}{}{}}
\LANGpl{
\Question{
Rozważ wartości 
\( \log_{0{,}2}3;\) \(\log_20{,}3;\)  \(\log_{0{,}3}1;\) \(\log_23;\) \(\log_{\frac32}\frac23;\) \(\log_{\frac23}2;\) \(\log_{0{,}3}\frac32. \)
 Określ ile z podanych wartości jest ujemnych? Użyj wykresu funkcji logarytmicznej podanego poniżej.}
{\( 5 \)}{\( 4 \)}{\( 3 \)}{\( 2 \)}{}{}}
\LANGcs{
\Question{
Mějme hodnoty
\( \log_{0{,}2}3;\) \(\log_20{,}3;\)  \(\log_{0{,}3}1;\) \(\log_23;\) \(\log_{\frac32}\frac23;\) \(\log_{\frac23}2;\) \(\log_{0{,}3}\frac32. \)
 Určete kolik z daných hodnot je záporných. Použijte daný graf logaritmické funkce.}
{\( 5 \)}{\( 4 \)}{\( 3 \)}{\( 2 \)}{}{}}
\LANGsk{
\Question{
Na obrázku je graf logaritmickej funkcie. Na základe obrázku rozhodnite, koľko z uvedených hodnôt je záporných.
\( \log_{0{,}2}3;\) \(\log_20{,}3;\)  \(\log_{0{,}3}1;\) \(\log_23;\) \(\log_{\frac32}\frac23;\) \(\log_{\frac23}2;\) \(\log_{0{,}3}\frac32. \)}
{\( 5 \)}{\( 4 \)}{\( 3 \)}{\( 2 \)}{}{}}
\LANGes{
\Question{
Dados los valores
\( \log_{0.2}3;\) \(\log_20.3;\)  \(\log_{0.3}1;\) \(\log_23;\) \(\log_{\frac32}\frac23;\) \(\log_{\frac23}2;\) \(\log_{0.3}\frac32.\) Determina cuántos de los valores dados son negativos. Usa la gráfica de la función logarítmica representada abajo.}
{\( 5 \)}{\( 4 \)}{\( 3 \)}{\( 2 \)}{}{}}
\End

\Data\ID{1103136503}
\Updated{2020-07-15 03:07:21}
\Level{B}
\SubArea{Logarithmic Functions}
\IMG{1103136503.pdf}
\LANGen{
\Question{
What are the values of the real variable \( x \), if \( \log_{0.4}x \leq 0 \)? Use the graph of \( f(x)=\log_{0.4}x \) given below.}
{\( x\in[1;\infty) \)}{\( x\in(0;1] \)}{\( x\in(-\infty;0] \)}{\( x\in[0;\infty) \)}{}{}}
\LANGpl{
\Question{
Jakie są wartości rzeczywistej zmiennej \( x \), jeśli \( \log_{0{,}4}x \leq 0 \)? Użyj wykresu \( f(x)=\log_{0{,}4}x \) podanego poniżej.}
{\( x\in\langle1;\infty) \)}{\( x\in(0;1\rangle \)}{\( x\in(-\infty;0\rangle \)}{\( x\in\langle0;\infty) \)}{}{}}
\LANGcs{
\Question{
Jaké jsou hodnoty reálné proměnné \( x \), jestliže \( \log_{0{,}4}x \leq 0 \)? Použijte graf \( f(x)=\log_{0{,}4}x \).}
{\( x\in\langle1;\infty) \)}{\( x\in(0;1\rangle \)}{\( x\in(-\infty;0\rangle \)}{\( x\in\langle0;\infty) \)}{}{}}
\LANGsk{
\Question{
Na obrázku je graf funkcie \( f(x)=\log_{0{,}4}x \). Aké sú hodnoty reálnej premennej \( x \), ak \( \log_{0{,}4}x \leq 0 \)?}
{\( x\in\langle1;\infty) \)}{\( x\in(0;1\rangle \)}{\( x\in(-\infty;0\rangle \)}{\( x\in\langle0;\infty) \)}{}{}}
\LANGes{
\Question{
¿Cuáles son los valores de la variable real \( x \), si \( \log_{0.4}x \leq 0 \)? Usa la gráfica de la función \( f(x)=\log_{0.4}x \) dada abajo.}
{\( x\in[1;\infty) \)}{\( x\in(0;1] \)}{\( x\in(-\infty;0] \)}{\( x\in[0;\infty) \)}{}{}}
\End

\Data\ID{1103136504}
\Updated{2020-07-15 03:07:21}
\Level{B}
\SubArea{Logarithmic Functions}
\IMG{1103136504.pdf}
\LANGen{
\Question{
What are the values of the real variable \( x \), if \( \log_7x < 0 \)? Use the graph of \( f(x)=\log_7x \) given below.}
{\( x\in(0;1) \)}{\( x\in[1;\infty) \)}{\( x\in(0;\infty) \)}{\( x\in[0;1] \)}{}{}}
\LANGpl{
\Question{
Jakie są wartości rzeczywistej zmiennej \( x \), jeśli \( \log_7x < 0 \)? Użyj wykresu  \( f(x)=\log_7x \) podanego poniżej.}
{\( x\in(0;1) \)}{\( x\in\langle1;\infty) \)}{\( x\in(0;\infty) \)}{\( x\in\langle0;1\rangle \)}{}{}}
\LANGcs{
\Question{
Jaké jsou hodnoty reálné proměnné \( x \), jestliže \( \log_7x < 0 \)? Použijte graf \( f(x)=\log_7x \).}
{\( x\in(0;1) \)}{\( x\in\langle1;\infty) \)}{\( x\in(0;\infty) \)}{\( x\in\langle0;1\rangle \)}{}{}}
\LANGsk{
\Question{
Na obrázku je graf funkcie \( f(x)=\log_7x \). Aké sú hodnoty reálnej premennej \( x \), ak \( \log_7x < 0 \)?}
{\( x\in(0;1) \)}{\( x\in\langle1;\infty) \)}{\( x\in(0;\infty) \)}{\( x\in\langle0;1\rangle \)}{}{}}
\LANGes{
\Question{
¿Cuáles son los valores de la variable real \( x \) si \( \log_7x < 0 \)? Usa la gráfica de la función \( f(x)=\log_7x \) dada abajo.}
{\( x\in(0;1) \)}{\( x\in[1;\infty) \)}{\( x\in(0;\infty) \)}{\( x\in[0;1] \)}{}{}}
\End

\Data\ID{1103136505}
\Updated{2020-07-15 03:07:21}
\Level{B}
\SubArea{Logarithmic Functions}
\IMG{1103136505.pdf}
\LANGen{
\Question{
What are the values of the real variable \( x \), if \( \log_{\frac14}x > \log_{\frac14}3 \)? Use the graph of   \( f(x)=\log_{\frac14}x \) given below.}
{\( x\in(0;3) \)}{\( x\in(-\infty;3) \)}{\( x\in(3;\infty) \)}{\( x\in(0;\infty) \)}{}{}}
\LANGpl{
\Question{
Jakie są wartości rzeczywistej zmiennej  \( x \), jeśli \( \log_{\frac14}x > \log_{\frac14}3 \)? Użyj wykresu \( f(x)=\log_{\frac14}x \) podanego poniżej.}
{\( x\in(0;3) \)}{\( x\in(-\infty;3) \)}{\( x\in(3;\infty) \)}{\( x\in(0;\infty) \)}{}{}}
\LANGcs{
\Question{
Jaké jsou hodnoty reálné proměnné \( x \), jestliže \( \log_{\frac14}x > \log_{\frac14}3 \)? Použijte graf  \( f(x)=\log_{\frac14}x \).}
{\( x\in(0;3) \)}{\( x\in(-\infty;3) \)}{\( x\in(3;\infty) \)}{\( x\in(0;\infty) \)}{}{}}
\LANGsk{
\Question{
Na obrázku je graf funkcie \( f(x)=\log_{\frac14}x \).  Aké sú hodnoty reálnej premennej \( x \), ak \( \log_{\frac14}x > \log_{\frac14}3 \)?}
{\( x\in(0;3) \)}{\( x\in(-\infty;3) \)}{\( x\in(3;\infty) \)}{\( x\in(0;\infty) \)}{}{}}
\LANGes{
\Question{
¿Cuáles son los valores de la variable real \( x \), si \( \log_{\frac14}x > \log_{\frac14}3 \)? Usa la gráfica de la función \( f(x)=\log_{\frac14}x \) dada abajo.}
{\( x\in(0;3) \)}{\( x\in(-\infty;3) \)}{\( x\in(3;\infty) \)}{\( x\in(0;\infty) \)}{}{}}
\End

\Data\ID{1103136506}
\Updated{2020-07-15 03:07:21}
\Level{B}
\SubArea{Logarithmic Functions}
\IMG{1103136506.pdf}
\LANGen{
\Question{
What are the values of the real variable \( x \), if  \( \log_4x \leq \log_40.7 \)? Use the graph of \( f(x)=\log_4x \) given below.}
{\( x\in(0;0.7] \)}{\( x\in[0.7;\infty) \)}{\( x\in[1;\infty) \)}{\( x\in(-\infty;0.7] \)}{}{}}
\LANGpl{
\Question{
Jakie są wartości rzeczywistej zmiennej \( x \), jeśli  \( \log_4x \leq \log_40{,}7 \)? Użyj wykresu \( f(x)=\log_4x \) podanego poniżej.}
{\( x\in(0;0{,}7\rangle \)}{\( x\in\langle0{,}7;\infty) \)}{\( x\in\langle1;\infty) \)}{\( x\in(-\infty;0{,}7\rangle \)}{}{}}
\LANGcs{
\Question{
Jaké jsou hodnoty reálné proměnné \( x \), jestliže \( \log_4x \leq \log_40{,}7 \)? Použijte graf \( f(x)=\log_4x \).}
{\( x\in(0;0{,}7\rangle \)}{\( x\in\langle0{,}7;\infty) \)}{\( x\in\langle1;\infty) \)}{\( x\in(-\infty;0{,}7\rangle \)}{}{}}
\LANGsk{
\Question{
Na obrázku je graf funkcie \( f(x)=\log_4x \). Aké sú hodnoty reálnej premennej \( x \), ak \( \log_4x \leq \log_40{,}7 \)?}
{\( x\in(0;0{,}7\rangle \)}{\( x\in\langle0{,}7;\infty) \)}{\( x\in\langle1;\infty) \)}{\( x\in(-\infty;0{,}7\rangle \)}{}{}}
\LANGes{
\Question{
¿Cuáles son los valores de la variable real \( x \), si \( \log_4x \leq \log_40.7 \)? Usa la gráfica de la función \( f(x)=\log_4x \) dada abajo.}
{\( x\in(0;0.7] \)}{\( x\in[0.7;\infty) \)}{\( x\in[1;\infty) \)}{\( x\in(-\infty;0.7] \)}{}{}}
\End

\Data\ID{1103136510}
\Updated{2022-04-23 05:04:42}
\Level{B}
\SubArea{Logarithmic Functions}
\IMG{1103136510.pdf}
\LANGen{
\Question{
Identify which of the following relations is correct.  Use the graph of \( f(x)=\log_{0.5}x \) given below.}
{\( \log_{0.5}\frac13 > \log_{0.5}1 >\log_{0.5} \frac43 > \log_{0.5}2 > \log_{0.5}5 \)}{\( \log_{0.5}5 > \log_{0.5}2 > \log_{0.5} \frac43 > \log_{0.5}1 > \log_{0.5}\frac13 \)}{\( \log_{0.5}1 > \log_{0.5}\frac13 > \log_{0.5}\frac43 > \log_{0.5}2 > \log_{0.5}5 \)}{\( \log_{0.5}\frac13 > \log_{0.5}\frac43 > \log_{0.5}1 > \log_{0.5}2 > \log_{0.5}5 \)}{}{}}
\LANGpl{
\Question{
Określ, która z poniższych zależności jest poprawna.  Użyj wykresu \( f(x)=\log_{0{,}5}x \) podanego poniżej.}
{\( \log_{0{,}5}\frac13 > \log_{0{,}5}1 >\log_{0{,}5} \frac43 > \log_{0{,}5}2 > \log_{0{,}5}5 \)}{\( \log_{0{,}5}5 > \log_{0{,}5}2 > \log_{0{,}5} \frac43 > \log_{0{,}5}1 > \log_{0{,}5}\frac13 \)}{\( \log_{0{,}5}1 > \log_{0{,}5}\frac13 > \log_{0{,}5}\frac43 > \log_{0{,}5}2 > \log_{0{,}5}5 \)}{\( \log_{0{,}5}\frac13 > \log_{0{,}5}\frac43 > \log_{0{,}5}1 > \log_{0{,}5}2 > \log_{0{,}5}5 \)}{}{}}
\LANGcs{
\Question{
Určete, který z následujících vztahů je správný. Použijte daný graf  \( f(x)=\log_{0{,}5}x \).}
{\( \log_{0{,}5}\frac13 > \log_{0{,}5}1 >\log_{0{,}5} \frac43 > \log_{0{,}5}2 > \log_{0{,}5}5 \)}{\( \log_{0{,}5}5 > \log_{0{,}5}2 > \log_{0{,}5} \frac43 > \log_{0{,}5}1 > \log_{0{,}5}\frac13 \)}{\( \log_{0{,}5}1 > \log_{0{,}5}\frac13 > \log_{0{,}5}\frac43 > \log_{0{,}5}2 > \log_{0{,}5}5 \)}{\( \log_{0{,}5}\frac13 > \log_{0{,}5}\frac43 > \log_{0{,}5}1 > \log_{0{,}5}2 > \log_{0{,}5}5 \)}{}{}}
\LANGsk{
\Question{
Určte, ktorý z nasledujúcich vzťahov je správny.  Použite graf funkcie \( f(x)=\log_{0{,}5}x \) danej na obrázku.}
{\( \log_{0{,}5}\frac13 > \log_{0{,}5}1 >\log_{0{,}5} \frac43 > \log_{0{,}5}2 > \log_{0{,}5}5 \)}{\( \log_{0{,}5}5 > \log_{0{,}5}2 > \log_{0{,}5} \frac43 > \log_{0{,}5}1 > \log_{0{,}5}\frac13 \)}{\( \log_{0{,}5}1 > \log_{0{,}5}\frac13 > \log_{0{,}5}\frac43 > \log_{0{,}5}2 > \log_{0{,}5}5 \)}{\( \log_{0{,}5}\frac13 > \log_{0{,}5}\frac43 > \log_{0{,}5}1 > \log_{0{,}5}2 > \log_{0{,}5}5 \)}{}{}}
\LANGes{
\Question{
Identifica cuál de las siguientes relaciones es correcta. Usa la gráfica de la función \( f(x)=\log_{0.5}x \) dada abajo.}
{\( \log_{0.5}\frac13 > \log_{0.5}1 >\log_{0.5} \frac43 > \log_{0.5}2 > \log_{0.5}5 \)}{\( \log_{0.5}5 > \log_{0.5}2 > \log_{0.5} \frac43 > \log_{0.5}1 > \log_{0.5}\frac13 \)}{\( \log_{0.5}1 > \log_{0.5}\frac13 > \log_{0.5}\frac43 > \log_{0.5}2 > \log_{0.5}5 \)}{\( \log_{0.5}\frac13 > \log_{0.5}\frac43 > \log_{0.5}1 > \log_{0.5}2 > \log_{0.5}5 \)}{}{}}
\End

\Data\ID{2000014101}
\Updated{2022-11-02 09:11:42}
\Level{B}
\SubArea{Logarithmic Functions}
\LANGen{
\Question{
Find the domain of the function \(f(x)=\log_{2015}\left(\log_{\frac{1}{2015}}(\log_{2015}x)\right)\).}
{\((1;2015)\)}{\((2015;\infty)\)}{\((0;\infty)\)}{\((0;2015)\)}{}{}}
\LANGpl{
\Question{
Znajdź dziedzinę funkcji \(f(x)=\log_{2015}\left(\log_{\frac{1}{2015}}(\log_{2015}x)\right)\).}
{\((1;2015)\)}{\((2015;\infty)\)}{\((0;\infty)\)}{\((0;2015)\)}{}{}}
\LANGcs{
\Question{
Určete definiční obor funkce \(f(x)=\log_{2015}\left(\log_{\frac{1}{2015}}(\log_{2015}x)\right)\).}
{\((1;2015)\)}{\((2015;\infty)\)}{\((0;\infty)\)}{\((0;2015)\)}{}{}}
\LANGsk{
\Question{
Definičným oborom funkcie \(f(x)=\log_{2015}\left(\log_{\frac{1}{2015}}(\log_{2015}x)\right)\) je:}
{\((1;2015)\)}{\((2015;\infty)\)}{\((0;\infty)\)}{\((0;2015)\)}{}{}}
\LANGes{
\Question{
Halla el dominio de la función \(f(x)=\log_{2015}\left(\log_{\frac{1}{2015}}(\log_{2015}x)\right)\).}
{\((1;2015)\)}{\((2015;\infty)\)}{\((0;\infty)\)}{\((0;2015)\)}{}{}}
\End

\Data\ID{2000014102}
\Updated{2022-11-02 09:11:42}
\Level{B}
\SubArea{Logarithmic Functions}
\LANGen{
\Question{
Make a true statement: The number \((\log_63)^2+(\log_62)^2+\log_64\cdot \log_63\) is}
{positive.}{smaller than 1.}{negative.}{irrational.}{}{}}
\LANGpl{
\Question{
Wskaż prawdziwą odpowiedź. Liczba \((\log_63)^2+(\log_62)^2+\log_64\cdot \log_63\) jest}
{dodatnia.}{mniejsza niż 1.}{ujemna.}{niewymierna.}{}{}}
\LANGcs{
\Question{
Doplňte pravdivé tvrzení: Číslo \((\log_63)^2+(\log_62)^2+\log_64\cdot \log_63\) je}
{kladné.}{menší než 1.}{záporné.}{iracionální.}{}{}}
\LANGsk{
\Question{
Doplňte pravdivé tvrdenie: Číslo \((\log_63)^2+(\log_62)^2+\log_64\cdot \log_63\) je}
{kladné.}{menšie ako jedna.}{záporné.}{iracionálne.}{}{}}
\LANGes{
\Question{
Completa correctamente la siguiente afirmación: El número \((\log_63)^2+(\log_62)^2+\log_64\cdot \log_63\) es}
{positivo.}{menor que 1.}{negativo.}{irracional.}{}{}}
\End

\Data\ID{2000014109}
\Updated{2022-11-02 09:11:46}
\Level{B}
\SubArea{Logarithmic Functions}
\LANGen{
\Question{
Identify which of the following relations is correct.}
{\( \log_3 10 >2\)}{\( \log_2 7 >3\)}{\( \log_2 3 < \log_3 2\)}{\( \log_4 15 >2\)}{}{}}
\LANGpl{
\Question{
Określ, która z poniższych relacji jest poprawna.}
{\( \log_3 10 >2\)}{\( \log_2 7 >3\)}{\( \log_2 3 < \log_3 2\)}{\( \log_4 15 >2\)}{}{}}
\LANGcs{
\Question{
Vyberte pravdivé tvrzení.}
{\( \log_3 10 >2\)}{\( \log_2 7 >3\)}{\( \log_2 3 < \log_3 2\)}{\( \log_4 15 >2\)}{}{}}
\LANGsk{
\Question{
Určte, ktoré z nasledujúcich tvrdení je pravdivé}
{\( \log_3 10 >2\)}{\( \log_2 7 >3\)}{\( \log_2 3 < \log_3 2\)}{\( \log_4 15 >2\)}{}{}}
\LANGes{
\Question{
Identifica cuál de las siguientes relaciones es correcta.}
{\( \log_3 10 >2\)}{\( \log_2 7 >3\)}{\( \log_2 3 < \log_3 2\)}{\( \log_4 15 >2\)}{}{}}
\End

\Data\ID{2010011009}
\Updated{2022-11-02 09:11:58}
\Level{B}
\SubArea{Logarithmic Functions}
\IMG{2010011001-figure0_0.pdf}
\LANGen{
\Question{
Identify which of the following relations is correct.  Use the graph of \( f(x)=\log_{\frac13}x \) given below.}
{\( \log_{\frac13}8 < \log_{\frac13}4< \log_{\frac13} 1 < \log_{\frac13}\frac12 < \log_{\frac13}\frac15 \)}{\( \log_{\frac13}\frac15 < \log_{\frac13}\frac12< \log_{\frac13} 1 < \log_{\frac13}4< \log_{\frac13}8 \)}{\( \log_{\frac13}\frac12 < \log_{\frac13}\frac15< \log_{\frac13} 1 < \log_{\frac13}4 < \log_{\frac13}8 \)}{\( \log_{\frac13}8 < \log_{\frac13}4< \log_{\frac13} 1 < \log_{\frac13}\frac15 < \log_{\frac13}\frac12 \)}{}{}}
\LANGpl{
\Question{
Określ, która z poniższych relacji jest poprawna.  Wykorzystaj podany poniżej wykres funkcji \( f(x)=\log_{\frac13}x \).}
{\( \log_{\frac13}8 < \log_{\frac13}4< \log_{\frac13} 1 < \log_{\frac13}\frac12 < \log_{\frac13}\frac15 \)}{\( \log_{\frac13}\frac15 < \log_{\frac13}\frac12< \log_{\frac13} 1 < \log_{\frac13}4< \log_{\frac13}8 \)}{\( \log_{\frac13}\frac12 < \log_{\frac13}\frac15< \log_{\frac13} 1 < \log_{\frac13}4 < \log_{\frac13}8 \)}{\( \log_{\frac13}8 < \log_{\frac13}4< \log_{\frac13} 1 < \log_{\frac13}\frac15 < \log_{\frac13}\frac12 \)}{}{}}
\LANGcs{
\Question{
Určete, která z daných relací je správná. Použijte graf \( f(x)=\log_{\frac13}x \) (viz obrázek).}
{\( \log_{\frac13}8 < \log_{\frac13}4< \log_{\frac13} 1 < \log_{\frac13}\frac12 < \log_{\frac13}\frac15 \)}{\( \log_{\frac13}\frac15 < \log_{\frac13}\frac12< \log_{\frac13} 1 < \log_{\frac13}4< \log_{\frac13}8 \)}{\( \log_{\frac13}\frac12 < \log_{\frac13}\frac15< \log_{\frac13} 1 < \log_{\frac13}4 < \log_{\frac13}8 \)}{\( \log_{\frac13}8 < \log_{\frac13}4< \log_{\frac13} 1 < \log_{\frac13}\frac15 < \log_{\frac13}\frac12 \)}{}{}}
\LANGsk{
\Question{
Určte, ktorý z nasledujúcich vzťahov je správny. Použite graf funkcie \( f(x)=\log_{\frac13}x \) danej na obrázku.}
{\( \log_{\frac13}8 < \log_{\frac13}4< \log_{\frac13} 1 < \log_{\frac13}\frac12 < \log_{\frac13}\frac15 \)}{\( \log_{\frac13}\frac15 < \log_{\frac13}\frac12< \log_{\frac13} 1 < \log_{\frac13}4< \log_{\frac13}8 \)}{\( \log_{\frac13}\frac12 < \log_{\frac13}\frac15< \log_{\frac13} 1 < \log_{\frac13}4 < \log_{\frac13}8 \)}{\( \log_{\frac13}8 < \log_{\frac13}4< \log_{\frac13} 1 < \log_{\frac13}\frac15 < \log_{\frac13}\frac12 \)}{}{}}
\LANGes{
\Question{
Identifica cuál de las siguientes relaciones es correcta.  Utiliza la gráfica de \( f(x)=\log_{\frac13}x \) dada a continuación.}
{\( \log_{\frac13}8 < \log_{\frac13}4< \log_{\frac13} 1 < \log_{\frac13}\frac12 < \log_{\frac13}\frac15 \)}{\( \log_{\frac13}\frac15 < \log_{\frac13}\frac12< \log_{\frac13} 1 < \log_{\frac13}4< \log_{\frac13}8 \)}{\( \log_{\frac13}\frac12 < \log_{\frac13}\frac15< \log_{\frac13} 1 < \log_{\frac13}4 < \log_{\frac13}8 \)}{\( \log_{\frac13}8 < \log_{\frac13}4< \log_{\frac13} 1 < \log_{\frac13}\frac15 < \log_{\frac13}\frac12 \)}{}{}}
\End

\Data\ID{2010016005}
\Updated{2022-11-02 09:11:14}
\Level{B}
\SubArea{Logarithmic Functions}
\LANGen{
\Question{
Let \(a=\log_3 \frac19\); \(b=\log_3 3\) and \(c=\log_3 \frac1{27}\). Which of the following statements is true?}
{\(c< a < b\)}{\(c < b < a\)}{\( b < c < a\)}{\( a < c < b\)}{}{}}
\LANGpl{
\Question{
Niech \(a=\log_3 \frac19\); \(b=\log_3 3\) i \(c=\log_3 \frac1{27}\). Które z poniższych stwierdzeń jest prawdziwe?}
{\(c< a < b\)}{\(c < b < a\)}{\( b < c < a\)}{\( a < c < b\)}{}{}}
\LANGcs{
\Question{
Nechť \(a=\log_3 \frac19\); \(b=\log_3 3\) a \(c=\log_3 \frac1{27}\). Který z následujících výroků je pravdivý?}
{\(c< a < b\)}{\(c < b < a\)}{\( b < c < a\)}{\( a < c < b\)}{}{}}
\LANGsk{
\Question{
Nech \(a=\log_3 \frac19\); \(b=\log_3 3\) a \(c=\log_3 \frac1{27}\). Ktoré z nasledujúcich tvrdení je pravdivé?}
{\(c< a < b\)}{\(c < b < a\)}{\( b < c < a\)}{\( a < c < b\)}{}{}}
\LANGes{
\Question{
Sean \(a=\log_3 \frac19\); \(b=\log_3 3\) y \(c=\log_3 \frac1{27}\). ¿Cuál de las siguientes afirmaciones es verdadera?}
{\(c< a < b\)}{\(c < b < a\)}{\( b < c < a\)}{\( a < c < b\)}{}{}}
\End

\Data\ID{2010016006}
\Updated{2022-11-02 09:11:14}
\Level{B}
\SubArea{Logarithmic Functions}
\LANGen{
\Question{
Let \(a=\log_4 \frac1{64}\); \(b=\log_4 4\) and \(c=\log_4 \frac1{16}\). Which of the following statements is true?}
{\(a< c < b\)}{\(b < c < a\)}{\( c < b < a\)}{\( a < b < c\)}{}{}}
\LANGpl{
\Question{
Niech \(a=\log_4 \frac1{64}\); \(b=\log_4 4\) i \(c=\log_4 \frac1{16}\). Które z poniższych stwierdzeń jest prawdziwe?}
{\(a< c < b\)}{\(b < c < a\)}{\( c < b < a\)}{\( a < b < c\)}{}{}}
\LANGcs{
\Question{
Nechť \(a=\log_4 \frac1{64}\); \(b=\log_4 4\) a \(c=\log_4 \frac1{16}\). Který z následujících výroků je pravdivý?}
{\(a< c < b\)}{\(b < c < a\)}{\( c < b < a\)}{\( a < b < c\)}{}{}}
\LANGsk{
\Question{
Nech \(a=\log_4 \frac1{64}\); \(b=\log_4 4\) a \(c=\log_4 \frac1{16}\). Ktoré z nasledujúcich tvrdení je pravdivé?}
{\(a< c < b\)}{\(b < c < a\)}{\( c < b < a\)}{\( a < b < c\)}{}{}}
\LANGes{
\Question{
Sean \(a=\log_4 \frac1{64}\); \(b=\log_4 4\) y \(c=\log_4 \frac1{16}\). ¿Cuál de las siguientes afirmaciones es verdadera?}
{\(a< c < b\)}{\(b < c < a\)}{\( c < b < a\)}{\( a < b < c\)}{}{}}
\End

\Data\ID{9000003803}
\Updated{2022-02-13 12:02:19}
\Level{B}
\SubArea{Logarithmic Functions}
\IMG{000038_03-figure0.pdf}
\LANGen{
\Question{
The function \(g\colon y =\log _{3}(x - 2)\) 
is graphed in the picture. In the following list identify a false statement.}
{The function \(g\) 
is a positive function.}{The domain of the function \(g\) 
is the interval \((2;\infty )\).}{The function \(g\) 
is not bounded.}{The function \(g\) 
is an increasing function.}{The function \(g\) 
has neither minimum nor maximum.}{The graph of the function \(g\) 
goes through \([5;1]\).}}
\LANGpl{
\Question{
Na rysunku poniżej przedstawiono wykres funkcji  \(g\colon y =\log _{3}(x - 2)\).
Które z poniższych zdań jest fałszywe?}
{Funkcja \(g\) 
jest funkcją dodatnią.}{Dziedzina funkcji \(g\) 
mieści się w przedziale \((2;\infty )\).}{Funkcja  \(g\) 
nie jest ograniczona.}{Funkcja  \(g\) 
jest funkcją rosnącą.}{Funkcja  \(g\) 
nie ma ani minimum ani maksimum.}{Wykres funkcji  \(g\) 
przechodzi przez punkt \([5;1]\).}}
\LANGcs{
\Question{
Je dána funkce \(g\colon y =\log _{3}(x - 2)\) 
(viz obrázek). Z následujících tvrzení vyberte to, které není pravdivé.}
{Funkce má všechny funkční hodnoty kladné.}{Definičním oborem funkce je interval
\((2;\infty )\).}{Funkce není omezená.}{Funkce je rostoucí.}{Funkce nenabývá maxima ani minima.}{Graf funkce \(g\) prochází bodem  \([5;1]\).}}
\LANGsk{
\Question{
Je daná funkcia \(g\colon y =\log _{3}(x - 2)\) 
na obrázku. Ktoré z nasledujúcich tvrdení nie je pravdivé?}
{Funkcia má všetky funkčné hodnoty kladné.}{Definičným oborom funkcie je interval
\((2;\infty )\).}{Funkcia nie je ohraničená.}{Funkcia je rastúca.}{Funkcia nemá maximum ani minimum.}{Graf funkcie \(g\) prechádza bodom \([5;1]\).}}
\LANGes{
\Question{
Dada la función \(g\colon y =\log _{3}(x - 2)\) (mira la imagen). Identifica el enunciado falso.}
{La función \(g\) 
es una función positiva.}{El dominio de la función \(g\) 
es el intervalo \((2;\infty )\).}{La función \(g\) 
no está acotada.}{La función \(g\) 
es una función creciente.}{La función \(g\) 
no tiene ni mínimos, ni máximos.}{La gráfica de la función \(g\) 
pasa por el punto \([5;1]\).}}
\End

\Data\ID{9000004808}
\Updated{2021-04-12 09:04:12}
\Level{B}
\SubArea{Logarithmic Functions}
\LANGen{
\Question{
In the following list identify a function which is bounded below.}
{\(y = 3^{x}\)}{\(y = -3^{x}\)}{\(y =\log _{3}x\)}{\(y = -\log _{3}x\)}{}{}}
\LANGpl{
\Question{
Która z podanych funkcji jest ograniczona z dołu?}
{\(y = 3^{x}\)}{\(y = -3^{x}\)}{\(y =\log _{3}x\)}{\(y = -\log _{3}x\)}{}{}}
\LANGcs{
\Question{
Která z následujících funkcí daných předpisem je omezená zdola?}
{\(y = 3^{x}\)}{\(y = -3^{x}\)}{\(y =\log _{3}x\)}{\(y = -\log _{3}x\)}{}{}}
\LANGsk{
\Question{
Z nasledujúcich predpisov vyberte zdola ohraničenú funkciu.}
{\(y = 3^{x}\)}{\(y = -3^{x}\)}{\(y =\log _{3}x\)}{\(y = -\log _{3}x\)}{}{}}
\LANGes{
\Question{
Identifica una función acotada inferiormente.}
{\(y = 3^{x}\)}{\(y = -3^{x}\)}{\(y =\log _{3}x\)}{\(y = -\log _{3}x\)}{}{}}
\End

\Data\ID{9000004810}
\Updated{2020-07-15 03:07:34}
\Level{B}
\SubArea{Logarithmic Functions}
\LANGen{
\Question{
In the following list identify a function which is not an increasing function.}
{\(y = 4x^{2}\)}{\(y =\log _{4}x\)}{\(y = 4x\)}{\(y = 4^{x}\)}{}{}}
\LANGpl{
\Question{
Która z niżej podanych funkcji nie jest funkcją rosnącą?}
{\(y = 4x^{2}\)}{\(y =\log _{4}x\)}{\(y = 4x\)}{\(y = 4^{x}\)}{}{}}
\LANGcs{
\Question{
Která z následujících funkcí daných předpisem není rostoucí na
svém definičním oboru?}
{\(y = 4x^{2}\)}{\(y =\log _{4}x\)}{\(y = 4x\)}{\(y = 4^{x}\)}{}{}}
\LANGsk{
\Question{
Ktorá z funkcií daných predpisom nie je rastúca na svojom definičnom obore?}
{\(y = 4x^{2}\)}{\(y =\log   _{4}x\)}{\(y = 4x\)}{\(y = 4^{x}\)}{}{}}
\LANGes{
\Question{
Identifica una función que no sea creciente.}
{\(y = 4x^{2}\)}{\(y =\log _{4}x\)}{\(y = 4x\)}{\(y = 4^{x}\)}{}{}}
\End

\Data\ID{9000004908}
\Updated{2020-07-15 03:07:34}
\Level{B}
\SubArea{Logarithmic Functions}
\LANGen{
\Question{
Complete the following statement: „The function
\(y =\log _{a^{2}-2a+2}x\) is
increasing if and only if ....”}
{\(a\in \mathbb{R}\setminus \{1\}\).}{\(a\in (-\infty ;\infty )\).}{\(a\in (0;\infty )\).}{\(a\in (1;\infty )\).}{}{}}
\LANGpl{
\Question{
Dokończ następujące zdanie: „Funkcja 
\(y =\log _{a^{2}-2a+2}x\) jest
 rosnąca wtedy i tylko wtedy gdy ....”}
{\(a\in \mathbb{R}\setminus \{1\}\).}{\(a\in (-\infty ;\infty )\).}{\(a\in (0;\infty )\).}{\(a\in (1;\infty )\).}{}{}}
\LANGcs{
\Question{
Funkce daná předpisem \(y =\log _{a^{2}-2a+2}x\) 
je rostoucí, právě když:}
{\(a\in \mathbb{R}\setminus \{1\}\)}{\(a\in (-\infty ;\infty )\)}{\(a\in (0;\infty )\)}{\(a\in (1;\infty )\)}{}{}}
\LANGsk{
\Question{
Funkcia daná predpisom \(y =\log _{a^{2}-2a+2}x\) 
je rastúca, práve vtedy keď:}
{\(a\in \mathbb{R}\setminus \{1\}\)}{\(a\in (-\infty ;\infty )\)}{\(a\in (0;\infty )\)}{\(a\in (1;\infty )\)}{}{}}
\LANGes{
\Question{
Completa el siguiente enunciado: "La función \(y =\log _{a^{2}-2a+2}x\) es creciente en el caso de que ..."}
{\(a\in \mathbb{R}\setminus \{1\}\).}{\(a\in (-\infty ;\infty )\).}{\(a\in (0;\infty )\).}{\(a\in (1;\infty )\).}{}{}}
\End

\Data\ID{1003101103}
\Updated{2021-04-12 08:04:43}
\Level{C}
\SubArea{Logarithmic Functions}
\LANGen{
\Question{
Find the true statement about the function  \( f(x)=\log_2|x| \).}
{The function \( f \) is even.}{The function \( f \) has the minimum at \( x=0 \).}{The function \( f \) is bounded.}{The function \( f \) is increasing.}{}{}}
\LANGpl{
\Question{
Które z podanych stwierdzeń dotyczących funkcji  \( f(x)=\log_2|x| \) jest prawdziwe?}
{Funkcja \( f \) jest parzysta.}{Funkcja   \( f \) osiąga minimum w \( x=0 \).}{Funkcja  \( f \) jest ograniczona.}{Funkcja \( f \) jest rosnąca.}{}{}}
\LANGcs{
\Question{
Funkce \( f \) je dána předpisem \( f(x)=\log_2|x| \). Vyberte pravdivý výrok.}
{Funkce \( f \) je sudá.}{Funkce \( f \) má minimum v bodě \( x=0 \).}{Funkce \( f \) je omezená.}{Funkce \( f \) je rostoucí.}{}{}}
\LANGsk{
\Question{
Vyberte pravdivé tvrdenie o funkcii \( f \) danej predpisom \( f(x)=\log_2|x| \).}
{Funkcia \( f \) párna.}{Funkcia \( f \) má minimum v bode \( x=0 \).}{Funkcia \( f \) je ohraničená.}{Funkcia \( f \) je rastúca.}{}{}}
\LANGes{
\Question{
Halla el enunciado que sea verdadero para la función \( f(x)=\log_2|x| \).}
{La función \( f \) es par.}{La función \( f \) tiene el mínimo en \( x=0 \).}{La función \( f \) está acotada.}{La función \( f \) es creciente.}{}{}}
\End

\Data\ID{1003138403}
\Updated{2020-07-15 03:07:21}
\Level{C}
\SubArea{Logarithmic Functions}
\LANGen{
\Question{
Find the domain of the function \( f(x)=\sqrt{\log_2x} \).}
{\( [ 1;\infty ) \)}{\( (0;\infty ) \)}{\( [ 0;\infty) \)}{\( (1;\infty) \)}{}{}}
\LANGpl{
\Question{
Wyznacz dziedzinę funkcji \( f(x)=\sqrt{\log_2x} \).}
{\( \langle 1;\infty ) \)}{\( (0;\infty ) \)}{\( \langle 0;\infty) \)}{\( (1;\infty) \)}{}{}}
\LANGcs{
\Question{
Určete definiční obor funkce \( f(x)=\sqrt{\log_2x} \).}
{\( \langle 1;\infty ) \)}{\( (0;\infty ) \)}{\( \langle 0;\infty) \)}{\( (1;\infty) \)}{}{}}
\LANGsk{
\Question{
Určte definičný obor funkcie \( f(x)=\sqrt{\log_2x} \).}
{\( \langle 1;\infty ) \)}{\( (0;\infty ) \)}{\( \langle 0;\infty) \)}{\( (1;\infty) \)}{}{}}
\LANGes{
\Question{
Halla el dominio de la función \( f(x)=\sqrt{\log_2x} \).}
{\( [ 1;\infty ) \)}{\( (0;\infty ) \)}{\( [ 0;\infty) \)}{\( (1;\infty) \)}{}{}}
\End

\Data\ID{1003138404}
\Updated{2020-07-15 03:07:21}
\Level{C}
\SubArea{Logarithmic Functions}
\LANGen{
\Question{
Find the range of the function \( f(x)=\left|\log_3x\right| \).}
{\( [ 0;\infty ) \)}{\( (0;\infty) \)}{\( (-\infty;\infty) \)}{\( [1;\infty) \)}{}{}}
\LANGpl{
\Question{
Wyznacz zakres funkcji \( f(x)=\left|\log_3x\right| \).}
{\( \langle 0;\infty ) \)}{\( (0;\infty) \)}{\( (-\infty;\infty) \)}{\( \langle1;\infty) \)}{}{}}
\LANGcs{
\Question{
Určete obor hodnot funkce \( f(x)=\left|\log_3x\right| \).}
{\( \langle 0;\infty ) \)}{\( (0;\infty) \)}{\( (-\infty;\infty) \)}{\( \langle1;\infty) \)}{}{}}
\LANGsk{
\Question{
Určte obor hodnôt funkcie \( f(x)=\left|\log_3x\right| \).}
{\( \langle 0;\infty ) \)}{\( (0;\infty) \)}{\( (-\infty;\infty) \)}{\( \langle1;\infty) \)}{}{}}
\LANGes{
\Question{
Halla el rango de la función \( f(x)=\left|\log_3x\right| \).}
{\( [ 0;\infty ) \)}{\( (0;\infty) \)}{\( (-\infty;\infty) \)}{\( [1;\infty) \)}{}{}}
\End

\Data\ID{1003170802}
\Updated{2021-04-12 09:04:01}
\Level{C}
\SubArea{Logarithmic Functions}
\LANGen{
\Question{
The acidity of a substance is measured by its \( pH \) value. In practice, \( pH \) is often assumed to be the negative logarithm of the hydrogen ion concentration \( H^+ \)  (i.e. the number of moles of hydrogen ions per liter of solution):
\[ pH = -\log{(H^+ )}. \]
If an apple juice has a \( pH \) of \( 3.8 \) and a grapefruit juice has a \( pH \) of \( 3.6 \) then the hydrogen ion concentration (\( \mathrm{mol}/\mathrm{l} \)) in the apple juice is how much less than in the grapefruit juice? Round the result to six decimal places.}
{\( 0.000093 \)}{\( -0.000093 \)}{\( 1.584893 \)}{\( 0 \)}{}{}}
\LANGpl{
\Question{
Kwasowość substancji mierzona jest za pomocą jej wartości \( pH \). W praktyce często przyjmuje się, że \( pH \)  jest logarytmem ujemnym stężenia jonów wodorowych \( H^+ \)  (t.j. liczba moli jonów wodoru na litr roztworu):
\[ pH = -\log{(H^+ )}. \]
Jeśli sok jabłkowy ma \( pH \) równe \( 3{,}8 \), a sok grejpfrutowy ma \( pH \) równe  \( 3{,}6 \) zatem o ile mniejsze jest stężenie jonów wodoru  (\( \mathrm{mol}/\mathrm{l} \)) w soku jabłkowym  w porównaniu do soku grejpfrutowego? Zaokrąglij do sześciu miejsc dziesiętnych.}
{\( 0{,}000093 \)}{\( -0{,}000093 \)}{\( 1{,}584893 \)}{\( 0 \)}{}{}}
\LANGcs{
\Question{
Kyselost látky se měří pomocí hodnoty \( pH \). V praxi se často \( pH \) vyjadřuje jako záporný logaritmus koncentrace vodíkových kationů \( H^+ \), tj. koncentrace vodíkových kationů se udává v molech na litr a vyjadřuje se rovnicí:
\[ pH = -\log{(H^+ )}. \]
Jablkový džus má hodnotu \( pH \) rovnou \( 3{,}8 \) a grepový džus má hodnotu \( pH \) rovnou \( 3{,}6 \). O kolik je koncentrace vodíkových kationů (\( \mathrm{mol}/\mathrm{l} \)) v jablkovém džuse menší než v grepovém? Výsledek zaokrouhlete na šest desetinných míst.}
{\( 0{,}000093 \)}{\( -0{,}000093 \)}{\( 1{,}584893 \)}{\( 0 \)}{}{}}
\LANGsk{
\Question{
Kyslosť látky sa meria pomocou hodnoty \( pH \). V praxi sa často \( pH \) vyjadruje ako záporný logaritmus koncentrácie vodíkových katiónov \( H^+ \), tj. koncentrácia vodíkových katiónov sa udáva v moloch na liter a vyjadruje sa rovnicou:
\[ pH = -\log{(H^+ )}. \]
Jablkový džús má hodnotu \( pH \) rovnú \( 3{,}8 \) a grepový džús má hodnotu \( pH \) rovnú \( 3{,}6 \). O koľko je koncentrácia vodíkových katiónov (\( \mathrm{mol}/\mathrm{l} \)) v jablkovom džúse menšia ako v grepovom? Výsledok zaokrúhlite na šesť desatinných miest.}
{\( 0{,}000093 \)}{\( -0{,}000093 \)}{\( 1{,}584893 \)}{\( 0 \)}{}{}}
\LANGes{
\Question{
El \( pH \) es la medida del grado de acidez de una sustancia. En la práctica, el \( pH \) es el logaritmo negativo de la concentración molar de los iones de hidrógeno \( H^+ \)  (es decir, el número de moles de iones de hidrógeno en un litro de una solución):
\[ pH = -\log{(H^+ )}. \]
Si un zumo de manzanas tiene el \( pH \) igual a \( 3.8 \) y un zumo de pomelo tiene el \( pH \) igual a \( 3.6 \), ¿Cuál es la diferencia de la concentración molar de los iones de hidrógeno (\( \mathrm{mol}/\mathrm{l} \)) entre el zumo de manzanas y el zumo de pomelo? Aproxima el resultado a seis cifras decimales.}
{\( 0.000093 \)}{\( -0.000093 \)}{\( 1.584893 \)}{\( 0 \)}{}{}}
\End

\Data\ID{1003170803}
\Updated{2020-07-15 03:07:21}
\Level{C}
\SubArea{Logarithmic Functions}
\LANGen{
\Question{
Loudness refers to how loud or soft a sound seems to a listener. The loudness \( L \) of sound in decibels, (\( \mathrm{dB} \)) is determined by the formula
\[ L=10\log{\left(\frac I{I_0}\right)}, \]
where \( I \) is intensity of the sound waves (in watts per square meter, \( \mathrm{W}/\mathrm{m}^2 \)) and \( \mathrm{I}_0 \) is fixed and determined intensity of the sound (\( 10^{-12}\,\mathrm{W}/\mathrm{m}^2 \))  that can barely be perceived. 
\[ \]
The intensity of the sound produced by the rock band is about \( 100 \) times greater than the intensity of the sound emitted by conventional motor vehicles. How much greater (in \( \mathrm{dB} \)) is loudness of rock band than loudness of motor vehicles?}
{\( 20 \)}{\( 100 \)}{\( 10 \)}{\( 2 \)}{}{}}
\LANGpl{
\Question{
Głośność odnosi się do tego jak głośny lub delikatny jest dźwięk dla słuchacza. Głośność \( L \)  dźwięku w decybelach (\( \mathrm{dB} \)) określona jest za pomocą wzoru 
\[ L=10\log{\left(\frac I{I_0}\right)}, \]
gdzie \( I \) oznacza intensywność fal dźwiękowych wyrażoną w Wattach na metr kwadratowy \( \mathrm{W}/\mathrm{m}^2 \)) i \( \mathrm{I}_0 \) jest stałe, a określona intensywność dźwięku (\( 10^{-12}\,\mathrm{W}/\mathrm{m}^2 \)), ledwie słyszalnego.
\[ \]
Intensywność dźwięku wytwarzanego przez zespół rockowy jest około  \( 100 \) razy większa niż intensywność dźwięku emitowanego przez standardowy silnik samochodowy. O ile wyższa (w \( \mathrm{dB} \)) jest głośność zespołu rockowego  od głośności silnika samochodowego?}
{\( 20 \)}{\( 100 \)}{\( 10 \)}{\( 2 \)}{}{}}
\LANGcs{
\Question{
Hladina intenzity zvuku \( L \) vyjadřuje úroveň „hlasitosti“ zvuku. Je měřená v decibelech (\( \mathrm{dB} \)) a vyjadřuje se vztahem
\[ L=10\log{\left(\frac I{I_0}\right)}, \]
kde \( I \) je intenzita zvuku (ve wattech na metr čtvereční \( \mathrm{W}/\mathrm{m}^2 \)) a \( \mathrm{I}_0 \) je prahová intenzita zvuku (tj. intenzita zodpovědná za slyšitelnost), která má hodnotu \( 10^{-12}\,\mathrm{W}/\mathrm{m}^2 \).
\[ \]
Intenzita zvuku vyprodukovaná rockovou kapelou je \( 100 \) krát větší než intenzita zvuku vyprodukovaná běžnými motorovými vozidly.
O kolik decibelů je hladina intenzity zvuku rockového koncertu větší oproti hlasitosti běžného motorového vozidla?}
{\( 20 \)}{\( 100 \)}{\( 10 \)}{\( 2 \)}{}{}}
\LANGsk{
\Question{
Hladina intenzity zvuku \( L \) vyjadruje úroveň „hlasitosti“ zvuku. Je meraná v decibeloch (\( \mathrm{dB} \)) a vyjadruje sa vzťahom
\[ L=10\log{\left(\frac I{I_0}\right)}, \]
kde \( I \) je intenzita zvuku (vo wattoch na meter štvorcový \( \mathrm{W}/\mathrm{m}^2 \)) a \( \mathrm{I}_0 \) je prahová intenzita zvuku (tj. intenzita zodpovedajúca počuteľnosti), ktorá má hodnotu \( 10^{-12}\,\mathrm{W}/\mathrm{m}^2 \).
\[ \]
Intenzita zvuku vyprodukovaná rockovou kapelou je \( 100 \) krát väčšia ako intenzita zvuku vyprodukovaná bežnými motorovými vozidlami.
O koľko decibelov je hladina intenzity zvuku rockového koncertu väčšia oproti hlasitosti bežného motorového vozidla?}
{\( 20 \)}{\( 100 \)}{\( 10 \)}{\( 2 \)}{}{}}
\LANGes{
\Question{
El nivel de la intensidad del sonido \( L \) se refiere al "volumen sonoro" (si es alto o bajo para el oyente) y se puede expresar en decibelios (\( \mathrm{dB} \)). La expresión es
\[ L=10\log{\left(\frac I{I_0}\right)}, \] 

donde \( I \) es la intensidad del sonido (en el vatio por metro cuadrado, \( \mathrm{W}/\mathrm{m}^2 \)) y \( \mathrm{I}_0 \) es el umbral de audición estándar (la intensidad de poder ser audible) determinado por (\( 10^{-12}\,\mathrm{W}/\mathrm{m}^2 \)).
\[ \]
La intensidad del sonido producida por una banda de rock es alrededor de \( 100 \) veces más alta que la intensidad del sonido emitida por los motores de vehículos normales. ¿Cuál es la diferencia (en \( \mathrm{dB} \)) entre la intensidad del sonido de la banda y de los motores?}
{\( 20 \)}{\( 100 \)}{\( 10 \)}{\( 2 \)}{}{}}
\End

\Data\ID{1103082705}
\Updated{2021-04-12 08:04:33}
\Level{C}
\SubArea{Logarithmic Functions}
\IMG{1103082705.pdf}
\LANGen{
\Question{
Function \( f \)  is given completely by the next graph. Identify which of the following statements is false.}
{\( f(x)=\log_2|x|;\ x\in[0.25;8] \)}{\( f(x)=|\log_2 x |;\ x\in[0.25;8] \)}{\( f(x)=|-\log_2 x|;\ x\in[0.25;8]  \)}{\( f(x)=\left|\log_{\frac12} x \right|;\ x\in[0.25;8] \)}{}{}}
\LANGpl{
\Question{
Funkcja \( f \) wyrażona jest za pomocą wykresu.  Określ, które z poniższych stwierdzeń jest  fałszywe.}
{\( f(x)=\log_2|x|;\ x\in\langle0{,}25;8\rangle \)}{\( f(x)=|\log_2 x |;\ x\in\langle0{,}25;8\rangle \)}{\( f(x)=|-\log_2 x|;\ x\in\langle0{,}25;8\rangle  \)}{\( f(x)=\left|\log_{\frac12} x \right|;\ x\in\langle0{,}25;8\rangle \)}{}{}}
\LANGcs{
\Question{
Funkce \( f \)  je dána následujícím grafem. Určete, které z následujících tvrzení je nepravdivé.}
{\( f(x)=\log_2|x|;\ x\in\langle0{,}25;8\rangle \)}{\( f(x)=|\log_2 x |;\ x\in\langle0{,}25;8\rangle \)}{\( f(x)=|-\log_2 x|;\ x\in\langle0{,}25;8\rangle  \)}{\( f(x)=\left|\log_{\frac12} x \right|;\ x\in\langle0{,}25;8\rangle \)}{}{}}
\LANGsk{
\Question{
Funkcia \( f \) je daná grafom. Určte, ktoré z následujúcich tvrdení o funkcii  \( f \)  je nepravdivé.}
{\( f(x)=\log_2|x|;\ x\in\langle0{,}25;8\rangle \)}{\( f(x)=|\log_2 x |;\ x\in\langle0{,}25;8\rangle \)}{\( f(x)=|-\log_2 x|;\ x\in\langle0{,}25;8\rangle  \)}{\( f(x)=\left|\log_{\frac12} x \right|;\ x\in\langle0{,}25;8\rangle \)}{}{}}
\LANGes{
\Question{
Dada la función \( f \)  mediante la gráfica de abajo. Identifica cuál de los siguientes enunciados es falso.}
{\( f(x)=\log_2|x|;\ x\in[0.25;8] \)}{\( f(x)=|\log_2 x |;\ x\in[0.25;8] \)}{\( f(x)=|-\log_2 x|;\ x\in[0.25;8]  \)}{\( f(x)=\left|\log_{\frac12} x \right|;\ x\in[0.25;8] \)}{}{}}
\End

\Data\ID{1103138401}
\Updated{2020-07-15 03:07:21}
\Level{C}
\SubArea{Logarithmic Functions}
\IMG{1103138401.pdf}
\LANGen{
\Question{
From the following functions choose the one which is graphed in the picture.}
{\( f(x)=\left|\log_{\frac12}x\right| \)}{\( g(x) = \log_{\frac12}|x| \)}{\( k(x)=\log_{\frac12}\sqrt x \)}{\( m(x)=\sqrt{\log_{\frac12}x} \)}{}{}}
\LANGpl{
\Question{
Z podanych funkcji wybierz tę, której wykres przedstawiono poniżej.}
{\( f(x)=\left|\log_{\frac12}x\right| \)}{\( g(x) = \log_{\frac12}|x| \)}{\( k(x)=\log_{\frac12}\sqrt x \)}{\( m(x)=\sqrt{\log_{\frac12}x} \)}{}{}}
\LANGcs{
\Question{
Určete předpis funkce, jejíž graf je znázorněný na obrázku.}
{\( f(x)=\left|\log_{\frac12}x\right| \)}{\( g(x) = \log_{\frac12}|x| \)}{\( k(x)=\log_{\frac12}\sqrt x \)}{\( m(x)=\sqrt{\log_{\frac12}x} \)}{}{}}
\LANGsk{
\Question{
Určte predpis funkcie, ktorej graf je znázornený na obrázku.}
{\( f(x)=\left|\log_{\frac12}x\right| \)}{\( g(x) = \log_{\frac12}|x| \)}{\( k(x)=\log_{\frac12}\sqrt x \)}{\( m(x)=\sqrt{\log_{\frac12}x} \)}{}{}}
\LANGes{
\Question{
De las siguientes funciones elige la cuál está representada en la imagen.}
{\( f(x)=\left|\log_{\frac12}x\right| \)}{\( g(x) = \log_{\frac12}|x| \)}{\( k(x)=\log_{\frac12}\sqrt x \)}{\( m(x)=\sqrt{\log_{\frac12}x} \)}{}{}}
\End

\Data\ID{1103138402}
\Updated{2020-07-15 03:07:21}
\Level{C}
\SubArea{Logarithmic Functions}
\IMG{1103138402.pdf}
\LANGen{
\Question{
From the following functions choose the one which is graphed in the picture.}
{\( k(x) = \log_2|x| \)}{\( m(x) = |\log_2 x| \)}{\( f(x) = \log_2\sqrt x \)}{\( g(x) = \sqrt{\log_2 x} \)}{}{}}
\LANGpl{
\Question{
Z podanych funkcji wybierz tę, którą przedstawiono za pomocą poniższego wykresu.}
{\( k(x) = \log_2|x| \)}{\( m(x) = |\log_2 x| \)}{\( f(x) = \log_2\sqrt x \)}{\( g(x) = \sqrt{\log_2 x} \)}{}{}}
\LANGcs{
\Question{
Určete předpis funkce, jejíž graf je znázorněný na obrázku.}
{\( k(x) = \log_2|x| \)}{\( m(x) = |\log_2 x| \)}{\( f(x) = \log_2\sqrt x \)}{\( g(x) = \sqrt{\log_2 x} \)}{}{}}
\LANGsk{
\Question{
Určte predpis funkcie, ktorej graf je znázornený na obrázku.}
{\( k(x) = \log_2|x| \)}{\( m(x) = |\log_2 x| \)}{\( f(x) = \log_2\sqrt x \)}{\( g(x) = \sqrt{\log_2 x} \)}{}{}}
\LANGes{
\Question{
De las siguientes funciones elige la cuál está representada en la imagen.}
{\( k(x) = \log_2|x| \)}{\( m(x) = |\log_2 x| \)}{\( f(x) = \log_2\sqrt x \)}{\( g(x) = \sqrt{\log_2 x} \)}{}{}}
\End

\Data\ID{1103170801}
\Updated{2021-04-12 09:04:51}
\Level{C}
\SubArea{Logarithmic Functions}
\LANGen{
\Question{
The acidity of a substance is measured by its \( pH \) value.  In practice, \( pH \) is often assumed to be the negative logarithm of the hydrogen ion concentration \( H^+ \)  (i.e. the number of moles of hydrogen ions per liter of solution):
\[ pH = -\log{(H^+ )}. \]
In one of the following pictures is a part of the \( pH-H^+ \)  graph. Choose the picture.}
{}{}{}{}{}{}}
\LANGpl{
\Question{
Kwasowość substancji mierzona jest na podstawie jej wartości \( pH \).  W praktyce,  \( pH \)
często przyjmuje się, że jest logarytmem ujemnym stężenia jonów wodorowych  \( H^+ \)  (t.j. liczba moli jonów wodoru na litr roztworu ):
\[ pH = -\log{(H^+ )}. \]
Na jednym z poniższych rysunków przedstawiono część wykresu \( pH-H^+ \). Wybierz ten rysunek}
{}{}{}{}{}{}}
\LANGcs{
\Question{
Kyselost látky se měří pomocí hodnoty \( pH \). V praxi se často \( pH \) vyjadřuje jako záporný logaritmus koncentrace vodíkových kationů \( H^+ \). Koncentrace vodíkových kationů se udává v molech na litr a vyjadřuje se rovnicí
\[ pH = -\log{(H^+ )}. \]
Který z následujících obrázků reprezentuje část grafu \( pH-H^+ \)?}
{}{}{}{}{}{}}
\LANGsk{
\Question{
Kyslosť látky sa meria pomocou hodnoty \( pH \). V praxi sa často \( pH \) vyjadruje ako záporný logaritmus koncentrácie vodíkových katiónov \( H^+ \). Koncentrácia vodíkových katiónov sa udáva v moloch na liter a vyjadruje sa rovnicou:
\[ pH = -\log{(H^+ )}. \]
Ktorý z nasledujúcich obrázkov reprezentuje časť grafu \( pH-H^+ \)?}
{}{}{}{}{}{}}
\LANGes{
\Question{
La acidez de sustancias se mide a través del \( pH \).  En la práctica, \( pH \) se define como el logaritmo negativo de la concentración de iones de hidróxido \( H^+ \)  (p.ej. la concentración de iones de hidróxido se mide en mol/L  y utiliza la siguiente ecuación):
\[ pH = -\log{(H^+ )}. \]
Elige la imagen que representa una parte de la gráfica de \( pH-H^+ \) .}
{}{}{}{}{}{}}

\IMGanswer{1103170801a_0.pdf}
\IMGanswer{1103170801b_0.pdf}
\IMGanswer{1103170801c_0.pdf}
\IMGanswer{1103170801d_0.pdf}\End

\Data\ID{2010011006}
\Updated{2022-11-02 09:11:54}
\Level{C}
\SubArea{Logarithmic Functions}
\LANGen{
\Question{
Find the domain of the following function.
\[ 
 f(x)= \log\left( \log_{\frac1{10}}\left(1-x^2\right)\right)
\]}
{\( (-1 ;0) \cup (0;1)\)}{\( (-1 ;1)\)}{\( (0;1)\)}{\([ 0;1)\)}{\((-\infty;1)\)}{}}
\LANGpl{
\Question{
Który z poniższych zbiorów jest dziedziną funkcji.
\[ 
 f(x)= \log\left( \log_{\frac1{10}}\left(1-x^2\right)\right)
\]}
{\( (-1 ;0) \cup (0;1)\)}{\( (-1 ;1)\)}{\( (0;1)\)}{\(\langle 0;1)\)}{\((-\infty;1)\)}{}}
\LANGcs{
\Question{
Určete definiční obor dané funkce.
\[ 
 f(x)= \log\left( \log_{\frac1{10}}\left(1-x^2\right)\right)
\]}
{\( (-1 ;0) \cup (0;1)\)}{\( (-1 ;1)\)}{\( (0;1)\)}{\(\langle 0;1)\)}{\((-\infty;1)\)}{}}
\LANGsk{
\Question{
Ktorá z uvedených množín je definičným oborom nasledujúci funkcie?\[ 
 f(x)= \log\left( \log_{\frac1{10}}\left(1-x^2\right)\right)\]}
{\( (-1 ;0) \cup (0;1)\)}{\( (-1 ;1)\)}{\( (0;1)\)}{\(\langle 0;1)\)}{\((-\infty;1)\)}{}}
\LANGes{
\Question{
Halla el dominio de la siguiente función.
\[ 
 f(x)= \log\left( \log_{\frac1{10}}\left(1-x^2\right)\right)
\]}
{\( (-1 ;0) \cup (0;1)\)}{\( (-1 ;1)\)}{\( (0;1)\)}{\([ 0;1)\)}{\((-\infty;1)\)}{}}
\End

\Data\ID{9000004905}
\Updated{2021-04-12 09:04:25}
\Level{C}
\SubArea{Logarithmic Functions}
\LANGen{
\Question{
In the following list identify a statement which is not true for the function
\(f\colon y = |\log (x - 3) - 1|\).}
{The function \(f\) 
is increasing on the domain.}{The domain of the function \(f\) 
is \((3;\infty )\).}{All values of the function \(f\) 
are nonnegative.}{The function \(f\) does
not have a \(y\)-intercept.}{The \(x\)-intercept
of the function \(f\) 
is \(x = 13\).}{The function \(f\) 
is not one-to-one.}}
\LANGpl{
\Question{
Które z poniższych zdań nie jest prawdziwe dla funkcji
\(f\colon y = |\log (x - 3) - 1|\)?}
{Funkcja  \(f\) 
jest rosnąca w dziedzinie.}{Dziedziną funkcji  \(f\) 
jest \((3;\infty )\).}{Wszystkie wartości funkcji \(f\) 
są nieujemne.}{Funkcja  \(f\) nie przecina
osi  \(y\) układu współrzędnych.}{Punktem przecięcia wykresu  z osią \(x\) układu współrzędnych 
funkcji  \(f\) 
jest \(x = 13\).}{Funkcja  \(f\) 
nie jest jednoznaczna.}}
\LANGcs{
\Question{
Které tvrzení popisující vlastnost funkce
\(f\colon y = |\log (x - 3) - 1|\) je
nepravdivé?}
{Funkce je rostoucí na celém definičním oboru.}{Definičním oborem funkce je interval
\((3;\infty )\).}{Všechny funkční hodnoty jsou nezáporné.}{Graf funkce nemá průsečík s osou
\(y\).}{Graf funkce protíná osu \(x\) 
v bodě \(x = 13\).}{Funkce není prostá.}}
\LANGsk{
\Question{
Ktoré z daných tvrdení o funkcii
\(f\colon y = |\log (x - 3) - 1|\) nie je
pravdivé?}
{Funkcia je rastúca na celom definičnom obore.}{Definičným oborom funkcie je interval
\((3;\infty )\).}{Všetky funkčné hodnoty sú nezáporné.}{Graf funkcie nemá priesečník s osou
\(y\).}{Graf funkcie pretína os \(x\) 
v bode \(x = 13\).}{Funkcia nie je prostá.}}
\LANGes{
\Question{
En la siguiente lista de enunciados elige el enunciado falso para la función 
\(f\colon y = |\log (x - 3) - 1|\).}
{La función \(f\) 
es creciente en todo el dominio.}{El dominio de la función \(f\) 
es \((3;\infty )\).}{Todos los valores de la función \(f\) 
son no negativos.}{La función \(f\) no tiene intersección con el eje \(y\).}{La intersección de la función \(f\) con el eje \(x\)
es \(x = 13\).}{La función \(f\) 
no es inyectiva.}}
\End

\Data\ID{9000033705}
\Updated{2022-04-23 05:04:42}
\Level{C}
\SubArea{Logarithmic Functions}
\LANGen{
\Question{
Find the domain of the following function. 
\[ 
 f(x) = \sqrt{\log (x^{2 } + 2x + 1)}
\]}
{\(\left (-\infty ;-2] \cup [ 0;\infty \right )\)}{\(\mathbb{R}\setminus \left \{-1\right \}\)}{\(\left (-1;\infty \right )\)}{\(\left (-\infty ;-1\right )\cup \left (1;\infty \right )\)}{\(\left (-\infty ;0\right )\cup \left (2;\infty \right )\)}{}}
\LANGpl{
\Question{
Wyznacz dziedzinę następującej funkcji.
\[ 
 f\colon y = \sqrt{\log (x^{2 } + 2x + 1)}
\]}
{\(\left (-\infty ;-2] \cup [ 0;\infty \right )\)}{\(\mathbb{R}\setminus \left \{-1\right \}\)}{\(\left (-1;\infty \right )\)}{\(\left (-\infty ;-1\right )\cup \left (1;\infty \right )\)}{\(\left (-\infty ;0\right )\cup \left (2;\infty \right )\)}{}}
\LANGcs{
\Question{
Definičním oborem funkce \(f\colon y = \sqrt{\log (x^{2 } + 2x + 1)}\) je množina:}
{\(\left (-\infty ;-2\rangle \cup \langle 0;\infty \right )\)}{\(\mathbb{R}\setminus \left \{-1\right \}\)}{\(\left (-1;\infty \right )\)}{\(\left (-\infty ;-1\right )\cup \left (1;\infty \right )\)}{\(\left (-\infty ;0\right )\cup \left (2;\infty \right )\)}{}}
\LANGsk{
\Question{
Definičným oborom funkcie \(f\colon y = \sqrt{\log (x^{2 } + 2x + 1)}\) je:}
{\(\left (-\infty ;-2\rangle \cup \langle 0;\infty \right )\)}{\(\mathbb{R}\setminus \left \{-1\right \}\)}{\(\left (-1;\infty \right )\)}{\(\left (-\infty ;-1\right )\cup \left (1;\infty \right )\)}{\(\left (-\infty ;0\right )\cup \left (2;\infty \right )\)}{}}
\LANGes{
\Question{
Halla el dominio de la siguiente función. 
\[ 
 f\colon y = \sqrt{\log (x^{2 } + 2x + 1)}
\]}
{\(\left (-\infty ;-2] \cup [ 0;\infty \right )\)}{\(\mathbb{R}\setminus \left \{-1\right \}\)}{\(\left (-1;\infty \right )\)}{\(\left (-\infty ;-1\right )\cup \left (1;\infty \right )\)}{\(\left (-\infty ;0\right )\cup \left (2;\infty \right )\)}{}}
\End
